% Group 44 report. Build with:  .venv/bin/python report/build.py --include-f5   (writes report/generated/* and main.pdf)
% Every result is a macro from report/generated/numbers.tex; do not type numbers here. Model-dependent sentences are
% macros too, so the text stays correct when new runs arrive. Figures that do not exist yet are omitted silently.
\documentclass[conference]{IEEEtran}
\IEEEoverridecommandlockouts
\usepackage[T1]{fontenc}
\usepackage{graphicx}
\usepackage{amsmath,amssymb}
\usepackage{booktabs,array,makecell}
\usepackage[table]{xcolor}
\usepackage{tcolorbox}
\usepackage{cite}
\usepackage{url}
\usepackage[hidelinks]{hyperref}

\definecolor{PendingColor}{HTML}{C0392B}
\definecolor{TooUniform}{HTML}{56B4E9}
\definecolor{TooContrasting}{HTML}{E69F00}
% In the submitted version a value whose input is missing prints as a neutral italic note.
\newcommand{\pending}[1]{\textit{#1}}
% Generated by report/build.py -- do not edit. Every value is computed from the pipeline outputs;
% \pending{...} marks values whose inputs do not exist yet.
\providecommand{\pending}[1]{{\color{PendingColor}\sffamily\bfseries\footnotesize[pending:~#1]}}
\newif\iffigfive
\figfivetrue
\newif\ifdraftstatus
\draftstatusfalse
\newif\iffigfifteen
\figfifteentrue
\newif\iffigfifteenb
\figfifteenbtrue
\newif\iffigfifteenc
\figfifteenctrue
\newcommand{\NAccRhoLoose}{0.70}
\newcommand{\NAccRhoStrict}{0.90}
\newcommand{\NAhnNames}{Cayenne, Yeast, Tequila and Sauerkraut}
\newcommand{\NAliasAccLinesPct}{97.9\%}
\newcommand{\NAliasAccRowsPct}{99.0\%}
\newcommand{\NAliasChosen}{tie-break preferences learned from expert lines}
\newcommand{\NAliasCoveragePct}{99.98\%}
\newcommand{\NAliasFullRange}{93.6--96.4\%}
\newcommand{\NAliasMatchedRange}{96.1--97.6\%}
\newcommand{\NAliasPartialPct}{2.3\%}
\newcommand{\NAliasSourceRange}{90.0--98.9\%}
\newcommand{\NAliasTestLines}{31{,}927}
\newcommand{\NAliasTestShare}{20}
\newcommand{\NAliasWorstSource}{Tarla Dalal (90.0\%)}
\newcommand{\NAppTabs}{three}
\newcommand{\NAucCMClf}{0.56}
\newcommand{\NAucCMComposite}{0.89}
\newcommand{\NAucCMCooc}{0.95}
\newcommand{\NAucCMNs}{0.50}
\newcommand{\NAucNsRange}{0.50--0.56}
\newcommand{\NAucNullClf}{0.63}
\newcommand{\NAucNullComposite}{0.94}
\newcommand{\NAucNullCooc}{0.99}
\newcommand{\NAucNullNs}{0.54}
\newcommand{\NAucOtherClf}{0.90}
\newcommand{\NAucOtherComposite}{0.87}
\newcommand{\NAucOtherCooc}{0.78}
\newcommand{\NAucOtherNs}{0.56}
\newcommand{\NAuditModelNames}{GPT-6-Luna, Gemini 3.5 Flash-Lite, Qwen3.5-9B and Llama-3.1-8B}
\newcommand{\NAuditModels}{four}
\newcommand{\NBestRegionF}{Indian Subcontinent (0.88)}
\newcommand{\NBonAblationText}{the independent classifier's share was 0.58 with all three components, 0.49 without the classifier, 0.55 without co-occurrence and 0.62 without $N_s$ consistency, while prevalence JSD was 0.272, 0.271, 0.273 and 0.303}
\newcommand{\NBonClfHighRange}{0.47--0.61}
\newcommand{\NBonClfLowRange}{0.31--0.47}
\newcommand{\NBonClfMonotone}{rose monotonically with $N$ for every model}
\newcommand{\NBonClfText}{0.45 to 0.61 (GPT-6-Luna), 0.31 to 0.47 (Gemini 3.5 Flash-Lite) and 0.47 to 0.60 (Qwen3.5-9B)}
\newcommand{\NBonDivMonotone}{fell monotonically with $N$ for every model}
\newcommand{\NBonDivText}{0.69 to 0.60 (GPT-6-Luna), 0.76 to 0.66 (Gemini 3.5 Flash-Lite) and 0.76 to 0.63 (Qwen3.5-9B)}
\newcommand{\NBonFpText}{0.25 to 0.18 (GPT-6-Luna), 0.34 to 0.03 (Gemini 3.5 Flash-Lite) and 0.38 to 0.13 (Qwen3.5-9B)}
\newcommand{\NBonJSDText}{0.264 to 0.275 (GPT-6-Luna), 0.275 to 0.269 (Gemini 3.5 Flash-Lite) and 0.259 to 0.270 (Qwen3.5-9B)}
\newcommand{\NBonModels}{three}
\newcommand{\NBonNMax}{8}
\newcommand{\NBonNValues}{1, 2, 4, 8}
\newcommand{\NBonNearCopy}{0}
\newcommand{\NBonNoveltyText}{0.62 to 0.64 (GPT-6-Luna), 0.59 to 0.62 (Gemini 3.5 Flash-Lite) and 0.56 to 0.60 (Qwen3.5-9B)}
\newcommand{\NBonPoolMax}{160}
\newcommand{\NBonPoolMin}{149}
\newcommand{\NBonSaiText}{1.01 to 1.17 (GPT-6-Luna), 0.80 to 1.04 (Gemini 3.5 Flash-Lite) and 0.92 to 1.15 (Qwen3.5-9B)}
\newcommand{\NBonSelectedMax}{20}
\newcommand{\NBonSelectedMin}{18}
\newcommand{\NBonSignText}{0.41 to 0.41 (GPT-6-Luna), 0.50 to 0.50 (Gemini 3.5 Flash-Lite) and 0.36 to 0.45 (Qwen3.5-9B)}
\newcommand{\NBoot}{1{,}000}
\newcommand{\NBudget}{5}
\newcommand{\NCMDgapPadjMax}{0.46}
\newcommand{\NCMDgapRange}{\ensuremath{-}0.16 to \ensuremath{-}0.02}
\newcommand{\NCMDgapSupport}{two of the four models (Qwen3.5-9B and Llama-3.1-8B)}
\newcommand{\NCMFpPadjMax}{0.039}
\newcommand{\NCMFpRange}{0.23--0.42}
\newcommand{\NCMFpRho}{0.02}
\newcommand{\NCMFpSupport}{all four models}
\newcommand{\NCMFreqMAE}{0.046}
\newcommand{\NCMFullFpRho}{0.08}
\newcommand{\NCMFullFreqMAE}{0.037}
\newcommand{\NCMFullNearCopyPct}{0.0\%}
\newcommand{\NCMFullPrevJSD}{0.60}
\newcommand{\NCMFullSAI}{0.16}
\newcommand{\NCMFullSignErrorPct}{48\%}
\newcommand{\NCMFullSignMatchPct}{47\%}
\newcommand{\NCMInstances}{10}
\newcommand{\NCMJSDPadjMax}{0.008}
\newcommand{\NCMJSDRange}{\ensuremath{-}0.45 to \ensuremath{-}0.43}
\newcommand{\NCMJSDSupport}{all four models}
\newcommand{\NCMMutations}{4}
\newcommand{\NCMNearCopyPct}{0.0\%}
\newcommand{\NCMPrevJSD}{0.62}
\newcommand{\NCMRandFpRho}{0.01}
\newcommand{\NCMRandFreqMAE}{0.039}
\newcommand{\NCMRandNearCopyPct}{0.0\%}
\newcommand{\NCMRandPrevJSD}{0.73}
\newcommand{\NCMRandSAI}{0.07}
\newcommand{\NCMRandSignErrorPct}{44\%}
\newcommand{\NCMRandSignMatchPct}{43\%}
\newcommand{\NCMSAI}{0.17}
\newcommand{\NCMSaiPadjMax}{0.008}
\newcommand{\NCMSaiRange}{\ensuremath{-}1.92 to \ensuremath{-}1.78}
\newcommand{\NCMSaiSupport}{all four models}
\newcommand{\NCMSeedPool}{20}
\newcommand{\NCMSignErrorPct}{40\%}
\newcommand{\NCMSignMatchPct}{46\%}
\newcommand{\NCallsPerModel}{704}
\newcommand{\NCallsTotal}{2{,}368}
\newcommand{\NCatJSDRange}{0.034--0.048}
\newcommand{\NCdbCategories}{21}
\newcommand{\NCdbCompound}{103}
\newcommand{\NCdbEntities}{1{,}033}
\newcommand{\NCdbIngredients}{930}
\newcommand{\NCdbLabels}{26}
\newcommand{\NCdbMiscLabels}{four}
\newcommand{\NCdbMiscRecipes}{207}
\newcommand{\NCdbPhrases}{161{,}208}
\newcommand{\NCdbRecipes}{45{,}772}
\newcommand{\NCdbRows}{456{,}279}
\newcommand{\NChiTopK}{30}
\newcommand{\NClfFOne}{0.39}
\newcommand{\NClfName}{linear SVM}
\newcommand{\NClfTopOne}{0.51}
\newcommand{\NClfTopThree}{0.74}
\newcommand{\NConditions}{2}
\newcommand{\NConditionsWord}{two}
\newcommand{\NContrastingAmongSmallest}{five}
\newcommand{\NContrastingCount}{six}
\newcommand{\NContrastingRobust}{six}
\newcommand{\NContrastingRobustText}{all six}
\newcommand{\NContribTopOne}{22/22}
\newcommand{\NContribTopThree}{63/66}
\newcommand{\NDRandomRange}{\ensuremath{-}0.49 to 0.91}
\newcommand{\NDirConsistent}{for every model, a larger share of contrasting than of uniform regions moved toward uniform pairing}
\newcommand{\NDirInterpretation}{This suggests that LLMs compress the pairing spectrum toward its middle; the pattern was not pre-registered, and it is not significant after BH correction.}
\newcommand{\NDirText}{3/6 vs 2/16 (GPT-6-Luna), 3/3 vs 0/5 (Gemini 3.5 Flash-Lite), 4/6 vs 6/16 (Qwen3.5-9B) and 5/6 vs 8/16 (Llama-3.1-8B)}
\newcommand{\NDivRange}{0.46--0.60}
\newcommand{\NDivRangePct}{46--60\%}
\newcommand{\NDocsPages}{7}
\newcommand{\NDrvJaccardRange}{0.00--0.43}
\newcommand{\NDrvKoreaText}{Korea's top real contributor, Soybean Sauce, is first for GPT-6-Luna and Qwen3.5-9B and absent from the top five for Llama-3.1-8B (Gemini 3.5 Flash-Lite: no Korean recipes yet)}
\newcommand{\NDrvNegText}{The negative median for Gemini 3.5 Flash-Lite reflects that $\chi\%$ is oriented by each cuisine's own $Z$ sign, so a model whose pairing sign differs from the real cuisine's can have a negative rank correlation.}
\newcommand{\NDrvPerModelText}{0.06 and 0.00 (GPT-6-Luna), 0.00 and \ensuremath{-}0.67 (Gemini 3.5 Flash-Lite), 0.25 and 0.76 (Qwen3.5-9B) and 0.43 and 0.79 (Llama-3.1-8B)}
\newcommand{\NDrvRefJaccard}{0.67}
\newcommand{\NDrvRefReps}{20}
\newcommand{\NDrvRefRho}{0.97}
\newcommand{\NDrvRhoRange}{\ensuremath{-}0.67 to 0.79}
\newcommand{\NEceCal}{0.06}
\newcommand{\NEceRaw}{0.34}
\newcommand{\NEqContrastingTotal}{19}
\newcommand{\NEqCuisinesTotal}{74}
\newcommand{\NEqGridMax}{160}
\newcommand{\NEqGridMin}{60}
\newcommand{\NEqGridStep}{10}
\newcommand{\NEqPerModelText}{2 too uniform and 11 too contrasting of 22 (GPT-6-Luna), 3 too uniform and 3 too contrasting of 8 (Gemini 3.5 Flash-Lite), 3 too uniform and 3 too contrasting of 22 (Qwen3.5-9B) and 9 too uniform and 2 too contrasting of 22 (Llama-3.1-8B)}
\newcommand{\NEqReps}{500}
\newcommand{\NEqShiftRange}{\ensuremath{-}0.29 to 0.11}
\newcommand{\NEqTypicalTotal}{38}
\newcommand{\NEqUniformTotal}{17}
\newcommand{\NEvalClfFOne}{0.40}
\newcommand{\NExactFiveRange}{94.0--99.9\%}
\newcommand{\NFPCanadaPct}{55\%}
\newcommand{\NFPChinaPct}{61\%}
\newcommand{\NFPFrancePct}{41\%}
\newcommand{\NFPRatePct}{8.6\%}
\newcommand{\NFPSouthEastAsiaPct}{17\%}
\newcommand{\NFPUsaPct}{6\%}
\newcommand{\NFeedbackItems}{3}
\newcommand{\NFindingCM}{Every model was closer to the real cuisines than a copy-mutate evolution baseline on ingredient prevalence, fingerprints and signature-ingredient use.}
\newcommand{\NFindingFingerprint}{their category fingerprints agreed less with the real cuisines than resampling allows (median $\rho$ 0.28--0.37 vs 0.75 for genuine samples)}
\newcommand{\NFindingPairing}{were not pushed toward uniform pairing (Wilcoxon $p\geq 0.064$ for every model)}
\newcommand{\NFindingSAI}{did not over-use signature ingredients on average (median stereotype amplification index 0.97--1.04)}
\newcommand{\NFpRhoRange}{0.28--0.37}
\newcommand{\NFreqMAERange}{0.031--0.051}
\newcommand{\NGbdtValFOne}{0.42}
\newcommand{\NGeminiAudRegions}{8}
\newcommand{\NGeminiCalls}{256}
\newcommand{\NGeminiEqContrasting}{3}
\newcommand{\NGeminiEqUniform}{3}
\newcommand{\NGeminiRecipes}{1{,}259}
\newcommand{\NGeminiRegions}{8}
\newcommand{\NGeminiSignAgreePct}{38\%}
\newcommand{\NGeminiSignErr}{3}
\newcommand{\NHOneEqFisherPRange}{0.018--0.229}
\newcommand{\NHOneEqFisherPadjRange}{0.071--0.229}
\newcommand{\NHOneEqFisherSupport}{none of the four models}
\newcommand{\NHOneEqSummary}{Wilcoxon $p\geq 0.064$ for every model}
\newcommand{\NHOneEqWilcoxPmin}{$p\geq 0.064$}
\newcommand{\NHOneEqWilcoxSupport}{none of the four models}
\newcommand{\NHOneExpDirection}{LLM cuisines were less uniform than the full cuisines}
\newcommand{\NHOneExpShiftRange}{\ensuremath{-}0.49 to \ensuremath{-}0.10}
\newcommand{\NHOneExpSupport}{none of the four models}
\newcommand{\NHOneFisherPadjMin}{1.00}
\newcommand{\NHOneFisherSupport}{none of the four models}
\newcommand{\NHTwoBDivPadjMax}{0.004}
\newcommand{\NHTwoBDivSupport}{all four models}
\newcommand{\NHTwoBPadjMin}{0.64}
\newcommand{\NHTwoBSupport}{none of the four models}
\newcommand{\NHTwoCBelowText}{14 of 22 (GPT-6-Luna), 4 of 8 (Gemini 3.5 Flash-Lite), 20 of 22 (Qwen3.5-9B) and 19 of 22 (Llama-3.1-8B)}
\newcommand{\NHTwoCPadjMax}{0.004}
\newcommand{\NHTwoCSummary}{Wilcoxon, BH-adjusted $p\leq 0.004$ for every model}
\newcommand{\NHTwoCSupport}{all four models}
\newcommand{\NKaggleEngineText}{every Kaggle run used the Hugging Face transformers fallback after vLLM failed on the T4}
\newcommand{\NKaggleLlamaHours}{2.7}
\newcommand{\NKaggleQwenHours}{4.9}
\newcommand{\NKaggleQwenMinutes}{293}
\newcommand{\NKappa}{10}
\newcommand{\NKoreaRealTopFive}{Soybean Sauce, Beef, Vinegar, Cayenne and Egg}
\newcommand{\NKoreaTopOne}{Soybean Sauce}
\newcommand{\NKoreaTopOneChi}{12.0}
\newcommand{\NKoreaTopOnePub}{16.85}
\newcommand{\NKoreaTopThree}{Vinegar}
\newcommand{\NKoreaTopThreeChi}{3.4}
\newcommand{\NKoreaTopThreePub}{6.15}
\newcommand{\NKoreaTopTwo}{Beef}
\newcommand{\NKoreaTopTwoChi}{4.5}
\newcommand{\NKoreaTopTwoPub}{7.70}
\newcommand{\NLLMN}{160}
\newcommand{\NLibModules}{14}
\newcommand{\NLibTests}{99}
\newcommand{\NLinesPerRecipeRange}{9.2--10.4}
\newcommand{\NLlamaAudRegions}{22}
\newcommand{\NLlamaCalls}{704}
\newcommand{\NLlamaEqContrasting}{2}
\newcommand{\NLlamaEqUniform}{9}
\newcommand{\NLlamaRecipes}{3{,}403}
\newcommand{\NLlamaRegions}{22}
\newcommand{\NLlamaSignAgreePct}{82\%}
\newcommand{\NLlamaSignErr}{2}
\newcommand{\NLogregValFOneRange}{0.38}
\newcommand{\NLunaAudRegions}{22}
\newcommand{\NLunaCalls}{704}
\newcommand{\NLunaEqContrasting}{11}
\newcommand{\NLunaEqUniform}{2}
\newcommand{\NLunaRecipes}{3{,}517}
\newcommand{\NLunaRegions}{22}
\newcommand{\NLunaSignAgreePct}{50\%}
\newcommand{\NLunaSignErr}{8}
\newcommand{\NMatchCanadaPct}{3\%}
\newcommand{\NMatchChinaPct}{5\%}
\newcommand{\NMatchFrancePct}{14\%}
\newcommand{\NMatchedN}{290}
\newcommand{\NMatchedReps}{200}
\newcommand{\NMaxOther}{3{,}000}
\newcommand{\NMaxTokens}{6{,}144}
\newcommand{\NMeetsProposal}{The D1 targets of the proposal are met.}
\newcommand{\NMinistralAudRegions}{\pending{Ministral-3-8B: extension: generation in progress}}
\newcommand{\NMinistralCalls}{\pending{Ministral-3-8B: extension: generation in progress}}
\newcommand{\NMinistralEqContrasting}{\pending{Ministral-3-8B: extension: generation in progress}}
\newcommand{\NMinistralEqUniform}{\pending{Ministral-3-8B: extension: generation in progress}}
\newcommand{\NMinistralRecipes}{\pending{Ministral-3-8B: extension: generation in progress}}
\newcommand{\NMinistralRegions}{\pending{Ministral-3-8B: extension: generation in progress}}
\newcommand{\NMinistralSignAgreePct}{\pending{Ministral-3-8B: extension: generation in progress}}
\newcommand{\NMinistralSignErr}{\pending{Ministral-3-8B: extension: generation in progress}}
\newcommand{\NMissingAdditives}{7}
\newcommand{\NMissingAhn}{4}
\newcommand{\NMissingCoarse}{13}
\newcommand{\NMissingCoreText}{none}
\newcommand{\NMissingExtraText}{Qwen3.5-4B and Ministral-3-8B}
\newcommand{\NMissingProfiles}{24}
\newcommand{\NModelNamesComplete}{GPT-6-Luna, Qwen3.5-9B and Llama-3.1-8B}
\newcommand{\NModelNamesDone}{GPT-6-Luna, Gemini 3.5 Flash-Lite, Qwen3.5-9B and Llama-3.1-8B}
\newcommand{\NModelsComplete}{three}
\newcommand{\NModelsDone}{four}
\newcommand{\NModelsDoneNum}{4}
\newcommand{\NNRand}{100{,}000}
\newcommand{\NNearCopy}{0.8}
\newcommand{\NNearCopyRangePct}{10--18\%}
\newcommand{\NNetBbBelowText}{22 of 22 (GPT-6-Luna), 8 of 8 (Gemini 3.5 Flash-Lite), 20 of 22 (Qwen3.5-9B) and 22 of 22 (Llama-3.1-8B)}
\newcommand{\NNetBbEdgesBelowText}{21 of 22 (GPT-6-Luna), 8 of 8 (Gemini 3.5 Flash-Lite), 18 of 22 (Qwen3.5-9B) and 22 of 22 (Llama-3.1-8B)}
\newcommand{\NNetBbEdgesRange}{67--104}
\newcommand{\NNetBbEdgesRef}{160}
\newcommand{\NNetBbOverlapBelowText}{0 of 22 (GPT-6-Luna), 3 of 8 (Gemini 3.5 Flash-Lite), 14 of 22 (Qwen3.5-9B) and 6 of 22 (Llama-3.1-8B)}
\newcommand{\NNetBbOverlapRange}{85--94\%}
\newcommand{\NNetBbOverlapRef}{98\%}
\newcommand{\NNetBbRange}{0.08--0.12}
\newcommand{\NNetBbRef}{0.29}
\newcommand{\NNetDraws}{100}
\newcommand{\NNetFiftyBelowText}{20 of 22 (GPT-6-Luna), 7 of 8 (Gemini 3.5 Flash-Lite), 19 of 22 (Qwen3.5-9B) and 16 of 22 (Llama-3.1-8B)}
\newcommand{\NNetFiftyRange}{0.07--0.09}
\newcommand{\NNetFiftyRef}{0.22}
\newcommand{\NNetHundredBelowText}{20 of 22 (GPT-6-Luna), 8 of 8 (Gemini 3.5 Flash-Lite), 19 of 22 (Qwen3.5-9B) and 19 of 22 (Llama-3.1-8B)}
\newcommand{\NNetHundredRange}{0.11--0.13}
\newcommand{\NNetHundredRef}{0.30}
\newcommand{\NNetSparser}{sparser than}
\newcommand{\NNoveltyRange}{0.61--0.68}
\newcommand{\NNullCandidates}{20{,}000}
\newcommand{\NOpenAISpend}{0.27}
\newcommand{\NOpenFamiliesDone}{two}
\newcommand{\NOpenModels}{Qwen3.5-9B and Llama-3.1-8B}
\newcommand{\NPairingMax}{16{,}030}
\newcommand{\NPairingMaxRegion}{USA}
\newcommand{\NPairingMin}{296}
\newcommand{\NPairingMinRegion}{Korea}
\newcommand{\NPairingRecipesTotal}{45{,}262}
\newcommand{\NParseAnyRange}{99.7--100.0\%}
\newcommand{\NParseOkRange}{70.3--99.9\%}
\newcommand{\NPartialText}{Gemini 3.5 Flash-Lite has returned 256 of its 704 calls (8 of 22 regions; subset set by the free-tier daily request cap).}
\newcommand{\NPersonaDgapPadjMax}{0.38}
\newcommand{\NPersonaDgapPadjMin}{0.18}
\newcommand{\NPersonaDgapRange}{\ensuremath{-}0.162 to \ensuremath{-}0.065}
\newcommand{\NPersonaDgapSupport}{none of the four models}
\newcommand{\NPersonaFpPadjMax}{0.71}
\newcommand{\NPersonaFpPadjMin}{0.005}
\newcommand{\NPersonaFpRange}{0.006--0.504}
\newcommand{\NPersonaFpSupport}{one of the four models (Qwen3.5-9B)}
\newcommand{\NPersonaJSDPadjMax}{0.016}
\newcommand{\NPersonaJSDPadjMin}{\ensuremath{<}0.001}
\newcommand{\NPersonaJSDRange}{\ensuremath{-}0.092 to \ensuremath{-}0.034}
\newcommand{\NPersonaJSDSupport}{all four models}
\newcommand{\NPersonaPlaces}{8}
\newcommand{\NPersonaSaiPadjMax}{0.95}
\newcommand{\NPersonaSaiPadjMin}{0.15}
\newcommand{\NPersonaSaiRange}{\ensuremath{-}0.081 to 0.002}
\newcommand{\NPersonaSaiSupport}{none of the four models}
\newcommand{\NPilotCalls}{44}
\newcommand{\NPrevJSDRange}{0.16--0.20}
\newcommand{\NPromptHash}{468c7d7f4c79}
\newcommand{\NPropCaveat}{{}}
\newcommand{\NPropDgap}{+0.095}
\newcommand{\NPropDgapP}{0.074}
\newcommand{\NPropFp}{+0.036}
\newcommand{\NPropFpP}{0.50}
\newcommand{\NPropJSD}{+0.016}
\newcommand{\NPropJSDP}{0.007}
\newcommand{\NPropModels}{GPT-6-Luna and Gemini 3.5 Flash-Lite}
\newcommand{\NPropSai}{+0.020}
\newcommand{\NPropSaiP}{0.82}
\newcommand{\NPropVerdict}{no metric showed the proprietary models to be significantly more faithful}
\newcommand{\NProprietaryDone}{two}
\newcommand{\NQ}{0.05}
\newcommand{\NQwenAudRegions}{22}
\newcommand{\NQwenCalls}{704}
\newcommand{\NQwenEqContrasting}{3}
\newcommand{\NQwenEqUniform}{3}
\newcommand{\NQwenRecipes}{3{,}477}
\newcommand{\NQwenRegions}{22}
\newcommand{\NQwenSignAgreePct}{68\%}
\newcommand{\NQwenSignErr}{3}
\newcommand{\NQwenSmallAudRegions}{\pending{Qwen3.5-4B: extension: generation in progress}}
\newcommand{\NQwenSmallCalls}{\pending{Qwen3.5-4B: extension: generation in progress}}
\newcommand{\NQwenSmallEqContrasting}{\pending{Qwen3.5-4B: extension: generation in progress}}
\newcommand{\NQwenSmallEqUniform}{\pending{Qwen3.5-4B: extension: generation in progress}}
\newcommand{\NQwenSmallRecipes}{\pending{Qwen3.5-4B: extension: generation in progress}}
\newcommand{\NQwenSmallRegions}{\pending{Qwen3.5-4B: extension: generation in progress}}
\newcommand{\NQwenSmallSignAgreePct}{\pending{Qwen3.5-4B: extension: generation in progress}}
\newcommand{\NQwenSmallSignErr}{\pending{Qwen3.5-4B: extension: generation in progress}}
\newcommand{\NRecipesPerCall}{5}
\newcommand{\NRecipesPerCallWord}{five}
\newcommand{\NRecipesPerCondition}{80}
\newcommand{\NRecipesPerModel}{3{,}520}
\newcommand{\NRecipesPerRegionModel}{160}
\newcommand{\NRecipesTarget}{10{,}000}
\newcommand{\NRecipesTotal}{11{,}656}
\newcommand{\NRefCatJSD}{0.0024}
\newcommand{\NRefFpRho}{0.75}
\newcommand{\NRefFpSubsamples}{200}
\newcommand{\NRefFreqMAE}{0.0075}
\newcommand{\NRefNearCopyPct}{12.5\%}
\newcommand{\NRefPrevJSD}{0.031}
\newcommand{\NRefSAI}{1.00}
\newcommand{\NRefSAIBandHi}{1.10}
\newcommand{\NRefSAIBandLo}{0.91}
\newcommand{\NRefSamplesTotal}{22{,}000}
\newcommand{\NRefSignMatchPct}{78\%}
\newcommand{\NRefSubsamples}{1{,}000}
\newcommand{\NRegions}{22}
\newcommand{\NReplCategoryN}{22}
\newcommand{\NReplCategoryRho}{0.94}
\newcommand{\NReplCategorySigns}{22}
\newcommand{\NReplFreqCatN}{22}
\newcommand{\NReplFreqCatRho}{0.98}
\newcommand{\NReplFreqCatSigns}{21}
\newcommand{\NReplFrequencyN}{22}
\newcommand{\NReplFrequencyRho}{0.98}
\newcommand{\NReplFrequencySigns}{21}
\newcommand{\NReplRandomN}{22}
\newcommand{\NReplRandomRho}{0.98}
\newcommand{\NReplRandomSigns}{22}
\newcommand{\NReplRhoRange}{0.94--0.98}
\newcommand{\NReplicates}{16}
\newcommand{\NReviseClfEndRange}{0.55--0.58}
\newcommand{\NReviseClfStartRange}{0.45--0.47}
\newcommand{\NReviseClfText}{0.45 to 0.58 (GPT-6-Luna) and 0.47 to 0.55 (Qwen3.5-9B)}
\newcommand{\NReviseFpEndRange}{0.35--0.41}
\newcommand{\NReviseFpStartRange}{0.25--0.38}
\newcommand{\NReviseFpText}{0.25 to 0.35 (GPT-6-Luna) and 0.38 to 0.41 (Qwen3.5-9B)}
\newcommand{\NReviseIdentityPctRange}{77--81\%}
\newcommand{\NReviseIdentityRange}{0.77--0.81}
\newcommand{\NReviseJSDEndRange}{0.241--0.245}
\newcommand{\NReviseJSDStartRange}{0.259--0.264}
\newcommand{\NReviseJSDText}{0.264 to 0.241 (GPT-6-Luna) and 0.259 to 0.245 (Qwen3.5-9B)}
\newcommand{\NReviseModels}{GPT-6-Luna and Qwen3.5-9B}
\newcommand{\NReviseRoundsText}{two rounds for GPT-6-Luna and one round for Qwen3.5-9B}
\newcommand{\NReviseSignEndRange}{0.59--0.68}
\newcommand{\NReviseSignStartRange}{0.36--0.41}
\newcommand{\NReviseSignText}{0.41 to 0.59 (GPT-6-Luna) and 0.36 to 0.68 (Qwen3.5-9B)}
\newcommand{\NReviseText}{{}}
\newcommand{\NReviseVerb}{raised}
\newcommand{\NRoundsMax}{2}
\newcommand{\NSAIRange}{0.97--1.04}
\newcommand{\NScaleText}{\pending{this contrast is the planned extension with Qwen3.5-4B, whose generation is in progress}}
\newcommand{\NSignAgreeRange}{38--82\%}
\newcommand{\NSignErrContrastingTotal}{1}
\newcommand{\NSignErrExpectedText}{1.9 (GPT-6-Luna), 1.2 (Gemini 3.5 Flash-Lite), 1.9 (Qwen3.5-9B) and 1.9 (Llama-3.1-8B)}
\newcommand{\NSignErrRegionsTotal}{74}
\newcommand{\NSignErrText}{8 of 22 (GPT-6-Luna), 3 of 8 (Gemini 3.5 Flash-Lite), 3 of 22 (Qwen3.5-9B) and 2 of 22 (Llama-3.1-8B)}
\newcommand{\NSignErrTotal}{16}
\newcommand{\NSignatureK}{10}
\newcommand{\NSmallestK}{seven}
\newcommand{\NSplitTest}{20}
\newcommand{\NSplitTrain}{70}
\newcommand{\NSplitVal}{10}
\newcommand{\NSteerPreModels}{Qwen3.5-9B and Llama-3.1-8B}
\newcommand{\NSteerPreN}{1, 2, 4, 8, 16}
\newcommand{\NSteerPrePool}{1{,}600}
\newcommand{\NSteerPreRegions}{10}
\newcommand{\NSteerPreSelected}{100}
\newcommand{\NSummaryUTC}{2026-11-23 10:57 UTC}
\newcommand{\NSvmValFOne}{0.40}
\newcommand{\NTemp}{0.27}
\newcommand{\NTemperature}{0.7}
\newcommand{\NTopP}{0.95}
\newcommand{\NUnanimousText}{British Isles and France were too uniform, and Greece was too contrasting, for every model that covers them}
\newcommand{\NUniformCount}{16}
\newcommand{\NUniformFlipped}{Canada, China and France}
\newcommand{\NUniformFlippedCount}{three}
\newcommand{\NUniformFlippedShares}{16\%, 22\% and 28\%}
\newcommand{\NUniformMiddle}{Indian Subcontinent, Mexico, USA and South East Asia}
\newcommand{\NUniformMiddleRange}{72--94\%}
\newcommand{\NUniformRobust}{nine}
\newcommand{\NUsaFullDeltaNs}{12.77}
\newcommand{\NUsaFullPool}{12.81}
\newcommand{\NUsaOwnPool}{\ensuremath{-}0.08}
\newcommand{\NUsaPool}{237}
\newcommand{\NVocab}{666}
\newcommand{\NWeighting}{equal}
\newcommand{\NWorstRegionF}{Australia \& NZ (0.08)}
\newcommand{\NZItalyOurs}{249.3}
\newcommand{\NZItalyPub}{263}
\newcommand{\NZKoreaOurs}{\ensuremath{-}56.7}
\newcommand{\NZKoreaPub}{\ensuremath{-}32}
\newcommand{\NZScandinaviaOurs}{\ensuremath{-}155.0}
\newcommand{\NZScandinaviaPub}{\ensuremath{-}148}


\graphicspath{{../figures/}}
\newcommand{\figfile}[2]{\IfFileExists{../figures/#1}{\includegraphics[width=#2]{#1}}{}}
\newcommand{\iffigure}[2]{\IfFileExists{../figures/#1}{#2}{}}
\newcommand{\figcode}[1]{\renewcommand{\thefigure}{F#1}}
\newcommand{\Ns}{N_s}
\newcommand{\dNs}{\Delta N_s}
\setcellgapes{1pt}
\renewcommand{\topfraction}{0.92}
\renewcommand{\dbltopfraction}{0.92}
\renewcommand{\textfraction}{0.06}
\renewcommand{\floatpagefraction}{0.8}
\renewcommand{\dblfloatpagefraction}{0.8}
\setcounter{topnumber}{3}
\setcounter{dbltopnumber}{3}
\setcounter{totalnumber}{5}
\extrafloats{100}

\begin{document}

\title{Do AI Chefs Respect Culinary Fingerprints? Auditing and Steering LLM-Generated Recipes across 22 World Cuisines}

\author{\IEEEauthorblockN{Group 44 (roll nos.\ 2023355 and 2023477)}
\IEEEauthorblockA{BIO544 Computational Gastronomy (CGAS), IIIT-Delhi, Monsoon 2026\\
Instructor: Prof.\ Ganesh Bagler}}

\maketitle

\begin{abstract}
Large language models (LLMs) are routinely asked for ``traditional'' recipes, but whether their recipes carry the measurable culinary fingerprints of the cuisines they claim has not been tested. We built LLMCuisine-22: \NRecipesTotal{} recipes for the \NRegions{} CulinaryDB regions from \NModelsDone{} LLMs (\NModelNamesDone), requested with plain and home-cook-persona prompts and mapped to CulinaryDB entities with a reconstruction of the CulinaryDB aliasing protocol (\NAliasAccLinesPct{} agreement with held-out expert mappings). We audited them with the food-pairing framework of Singh and Bagler after reproducing its Fig.~4 (\NReplRandomSigns/\NReplRandomN{} signs, Spearman $\rho=\NReplRandomRho$). Calibrating the framework at the LLMs' sample size showed that the own-pool random null depends on the number of recipes: the pre-registered sign-error rule fires on \NFPRatePct{} of genuine \NLLMN-recipe samples, so we compare LLM cuisines with genuine samples of equal size. At equal size, LLM cuisines \NFindingPairing{} and \NFindingSAI, but \NFindingFingerprint, they used \NDivRangePct{} of the real cuisines' ingredient diversity, and they reproduced few of each cuisine's distinctive ingredient pairings (top-100 NPMI overlap \NNetHundredRange{} vs \NNetHundredRef{} for genuine samples). \NFindingCM{} A fidelity scorer (cuisine classifier, co-occurrence PMI, $\Ns$ consistency) separates held-out real recipes from other-cuisine, copy-mutate and null recipes (AUC \NAucOtherComposite, \NAucCMComposite, \NAucNullComposite). Best-of-$N$ reranking with it raised an independent classifier's agreement with the claimed cuisine from \NBonClfLowRange{} to \NBonClfHighRange{} but lowered diversity; critique-and-revise \NReviseVerb{} it from \NReviseClfStartRange{} to \NReviseClfEndRange{} while keeping \NReviseIdentityPctRange{} of each dish's ingredients. We release the dataset, the FlavorAudit library and the FlavorPrint app.
\end{abstract}

\ifdraftstatus
\begin{tcolorbox}[colback=yellow!5,colframe=PendingColor!60,boxrule=0.4pt,left=2pt,right=2pt,top=1pt,bottom=1pt,fontupper=\scriptsize]
\textbf{Draft status (generated by \texttt{report/build.py}).} Numbers come from \texttt{results/summary.json} (\NSummaryUTC). Items that change when their input exists:
\begin{itemize}\setlength{\itemsep}{0pt}
  \item Gemini 3.5 Flash-Lite: subset set by the free-tier daily request cap
  \item Qwen3.5-4B: extension: generation in progress
  \item Ministral-3-8B: extension: generation in progress
  \item this contrast is the planned extension with Qwen3.5-4B, whose generation is in progress
  \item Publishing D1 (Kaggle/HF/Zenodo), D4 (PyPI) and D5 (HF Space): team approval
\end{itemize}

\end{tcolorbox}
\fi

\begin{IEEEkeywords}
computational gastronomy, food pairing, culinary fingerprint, CulinaryDB, FlavorDB, large language models, recipe generation
\end{IEEEkeywords}

% =====================================================================================================
\section{Introduction}
People now ask LLMs for recipes as readily as they search for them, yet, as Bagler and Goel note, the text-generation abilities of language models ``with multi-sensorial consequences for taste and odor remain untested hitherto''~\cite{bagler2024}. Work on LLMs and cuisine so far measures text: cultural adaptation of recipes~\cite{cao2024}, LLM-judged authenticity in cuisine transfer~\cite{lee2024ash}, diversity in cross-cultural adaptation~\cite{hu2025carriage} and textual novelty against cultural distance, where LLMs failed to ground adaptations in culturally salient ingredients~\cite{carichon2026}. None of these looks at flavor chemistry.

Traditional cuisines have quantitative fingerprints. Recipes share flavor compounds between ingredient pairs more or less than chance (food pairing~\cite{ahn2011}); the \NRegions{} world regions of CulinaryDB split into \NUniformCount{} with uniform (positive) and \NContrastingCount{} with contrasting (negative) pairing, and ingredient popularity largely explains the pattern~\cite{singh2018}; a cuisine's category contributions form its ``culinary fingerprint''~\cite{jain2015}; relative prevalence measures authenticity~\cite{jain2018}; and copy-mutate models reproduce ingredient popularity~\cite{kinouchi2008,jain2018,tuwani2019}. These measures make an audit of LLM recipes possible. We ask four questions, pre-registered with their hypotheses before any full-scale LLM data were analysed (Section~\ref{sec:stats}):
\begin{itemize}
\item \textbf{RQ1} Do LLM cuisines reproduce each region's food-pairing sign and strength? (H1: LLMs push cuisines toward uniform pairing, so sign errors cluster in the contrasting regions.)
\item \textbf{RQ2} Do they rely on the same flavor drivers (RQ2a), amplify stereotypes and lose diversity (RQ2b; H2b: signature ingredients over-used, $\mathrm{SAI}>1$), and keep the category fingerprint and co-occurrence structure (RQ2c; H2c: fingerprints agree less with the real cuisine than genuine resamples do)?
\item \textbf{RQ3} Can fingerprint feedback steer generation toward fidelity without copying?
\item \textbf{RQ4} Do prompt framing, proprietary vs open models and scale matter? (H4: the persona helps slightly; proprietary or larger models are not consistently more faithful.)
\end{itemize}
\iffigure{F1_pipeline.pdf}{Fig.~\ref{fig:F1} outlines the pipeline. }The box ``What is new, and what is reused'' separates our contributions from the methods we reuse.

\begin{figure}[t]
\begin{tcolorbox}[colback=gray!4,colframe=gray!55,boxrule=0.4pt,left=3pt,right=3pt,top=2pt,bottom=2pt,fontupper=\small,title={What is new, and what is reused},fonttitle=\small\bfseries,colbacktitle=gray!20,coltitle=black]
\textbf{New in this project:} (i) LLMCuisine-22, an LLM recipe dataset aliased exactly as CulinaryDB was; (ii) the first audit of LLM recipes against flavor-compound pairing, culinary fingerprints and co-occurrence structure, with genuine equal-size samples and a copy-mutate model as references; (iii) the finding that the own-pool food-pairing null depends on sample size, with an equal-$N$ comparison that corrects for it; (iv) fingerprint-guided steering (best-of-$N$, critique-and-revise) evaluated only on held-out metrics; (v) the open FlavorAudit library and FlavorPrint app.\\[2pt]
\textbf{Reused methods (not contributions):} $\Ns$, the four null models, $Z$ and $\chi$~\cite{ahn2011,jain2015,singh2018}; $\Delta N_s^{\mathrm{cat}}$ fingerprints~\cite{jain2015}; relative prevalence~\cite{jain2018}; copy-mutate models~\cite{kinouchi2008,jain2018,tuwani2019}, also practised in our course assignments; NPMI~\cite{bouma2009} and the disparity filter~\cite{serrano2009}.
\end{tcolorbox}
\end{figure}

\IfFileExists{../figures/F1_pipeline.pdf}{%
\begin{figure*}[t]
\centering\figfile{F1_pipeline.pdf}{0.9\textwidth}
\figcode{1}\caption{How does the project go from real cuisines to audited and steered LLM recipes? Real CulinaryDB cuisines with FlavorDB profiles pass the replication gate; LLM recipes are generated, parsed and aliased with the CulinaryDB protocol, audited with the same functions as real and copy-mutate cuisines, and steered with the fidelity scorer. The bottom row lists the five deliverables (D1--D5).}\label{fig:F1}
\end{figure*}}{}

% =====================================================================================================
\section{Related Work}
\textbf{Food pairing and cuisine structure.} Ahn et al.\ introduced the flavor network and found positive food pairing in North American and Western European recipes and negative pairing in East Asian ones~\cite{ahn2011}. Jain, Rakhi and Bagler found negative pairing in Indian regional cuisines, defined $Z$ and the ingredient contribution $\chi$, and introduced category contributions as culinary fingerprints~\cite{jain2015}. Singh and Bagler compiled CulinaryDB, used four null models and identified uniform and contrasting cuisines and the ingredients that drive them~\cite{singh2018}. Flavor profiles come from FlavorDB~\cite{garg2018}, now served as FlavorDB2~\cite{goel2024flavordb2}. Copy-mutate models explain the rank-frequency structure of cuisines~\cite{kinouchi2008}; ranked (frequency) fitness fits Indian cuisines better~\cite{jain2018}, and Tuwani et al.\ applied the models to world cuisines~\cite{tuwani2019}. Cuisine classifiers from recipe content~\cite{sharma2020} motivate our scorer.

\textbf{LLMs and recipes.} Cao et al.\ adapt recipes across cultures~\cite{cao2024}, Lee et al.\ judge cuisine transfer for authenticity, sensitivity and harmony~\cite{lee2024ash}, Hu et al.\ use retrieval to diversify adaptations~\cite{hu2025carriage} and Carichon et al.\ show that LLM adaptations do not track cultural distance~\cite{carichon2026}. All evaluate text or judgments; none tests flavor-compound sharing or fingerprints. For steering we build on Self-Refine~\cite{madaan2023}; the dataset follows the datasheet template~\cite{gebru2021}.

% =====================================================================================================
\section{Dataset}
\subsection{Real recipes and flavor profiles}
\textbf{CulinaryDB}~\cite{singh2018} has \NCdbRecipes{} recipes under \NCdbLabels{} cuisine labels: the \NRegions{} analysis regions and \NCdbMiscLabels{} aggregate-only labels (\NCdbMiscRecipes{} recipes). Its \NCdbIngredients{} basic and \NCdbCompound{} compound ingredients fall into \NCdbCategories{} categories. Its \NCdbRows{} recipe--ingredient rows keep the raw line next to the curated entity, giving \NCdbPhrases{} distinct expert-mapped phrases, the gold standard for our aliaser. The pairing analysis uses the \NPairingRecipesTotal{} recipes of the \NRegions{} regions with at least two ingredients, from \NPairingMin{} (\NPairingMinRegion) to \NPairingMax{} (\NPairingMaxRegion).

\textbf{Flavor profiles.} We fetched FlavorDB~\cite{garg2018} on 2026-11-22; both public endpoints serve the same current data, and the 2018 version used by Singh and Bagler is no longer online. Following the paper, we reconstructed the \NMissingProfiles{} CulinaryDB ingredients without a FlavorDB entry: \NMissingCoarse{} coarse-grained items take their FlavorDB parent's profile, \NAhnNames{} take Ahn et al.'s compounds~\cite{ahn2011} mapped by CAS number to PubChem, and of the \NMissingAdditives{} additives MSG and citric acid get their own molecule while the others stay empty. Compound ingredients take the union of their constituents' profiles, giving profiles for all \NCdbEntities{} entities.

\subsection{LLMCuisine-22 (D1)}
Each model receives \NRegions{} regions $\times$ \NConditionsWord{} prompt conditions $\times$ \NReplicates{} replicates, each call asking for \NRecipesPerCallWord{} recipes: \NCallsPerModel{} calls and \NRecipesPerModel{} recipes per model (\NRecipesPerRegionModel{} per region, \NRecipesPerCondition{} per condition). The \emph{plain} prompt asks for traditional recipes of the region; the \emph{persona} prompt makes the model a home cook from one of \NPersonaPlaces{} places of the region, each used equally often (Appendix~\ref{app:prompts}). The prompts never mention flavor and were frozen (prompt set v1, sha256 \texttt{\NPromptHash\ldots}) after a \NPilotCalls-call pilot per model. Open-weight models use temperature \NTemperature{} and top-$p$ \NTopP{} with thinking off; API models use provider defaults at the lowest reasoning level; all use at most \NMaxTokens{} output tokens. Every call is logged (prompts, model id and revision, parameters, seed, usage, latency, raw output, parse status). Parsing repairs structure only, never content, and keeps the first \NRecipesPerCallWord{} recipes of a call. GPT-6-Luna ran through the OpenAI Batch API (total OpenAI spend \$\NOpenAISpend), Gemini 3.5 Flash-Lite on the free API, and the open-weight models on Kaggle T4 GPUs (\NKaggleEngineText; the full runs took \NKaggleQwenHours{} h for Qwen3.5-9B and \NKaggleLlamaHours{} h for Llama-3.1-8B). Llama-3.1-8B weights came from a mirror verified byte-identical to \texttt{meta-llama/Llama-3.1-8B-Instruct}; the Llama licence applies.

LLMCuisine-22 holds \NRecipesTotal{} recipes from \NCallsTotal{} calls (Table~\ref{tab:dataset}); \NMeetsProposal{} Between \NExactFiveRange{} of calls returned exactly five recipes, \NParseAnyRange{} parsed after structural repair, and \NAliasMatchedRange{} of ingredient lines were aliased. Table~\ref{tab:status} lists the runs that had not finished at submission; no result in this report depends on them. The release follows CulinaryDB's file format, with a datasheet~\cite{gebru2021}.

\textbf{Licences.} CulinaryDB and FlavorDB are CC BY-NC-SA 3.0, so the derived data are CC BY-NC-SA 4.0 with attribution; Ahn et al.'s data are CC BY 4.0. Code is MIT, and source data are downloaded on first use, never bundled in the package. Each LLM's output terms are recorded in the datasheet.

\begin{table*}[t]
\centering\caption{LLMCuisine-22 per model: calls (of the \NCallsPerModel{} planned), recipes, regions, instruction following, parsing and aliasing (full-mode calls only; first five recipes per call).}\label{tab:dataset}
\footnotesize\setlength{\tabcolsep}{4pt}% Generated by report/build.py -- do not edit.
\begin{tabular}{@{}lllrrrrrrrr@{}}
\toprule
Model & Family & Access & Calls & Recipes & Regions & \makecell{Exactly\\5 recipes} & \makecell{Parsed\\(strict)} & \makecell{Parsed after\\repair} & \makecell{Lines/\\recipe} & \makecell{Lines\\aliased} \\
\midrule
GPT-6-Luna & GPT (OpenAI) & OpenAI Batch API & 704/704 & 3{,}517 & 22 & 99.9\% & 99.9\% & 100.0\% & 10.4 & 96.1\% \\
Gemini 3.5 Flash-Lite & Gemini (Google) & Gemini API (free) & 256/704 & 1{,}259 & 8 & 94.5\% & 79.7\% & 100.0\% & 9.8 & 97.6\% \\
Qwen3.5-9B & Qwen (Alibaba) & Kaggle GPU & 704/704 & 3{,}477 & 22 & 98.2\% & 70.3\% & 99.7\% & 9.9 & 96.6\% \\
Llama-3.1-8B & Llama (Meta) & Kaggle GPU & 704/704 & 3{,}403 & 22 & 94.0\% & 93.2\% & 100.0\% & 9.2 & 97.3\% \\
Qwen3.5-4B & Qwen (Alibaba) & Kaggle GPU & \multicolumn{8}{l}{\pending{extension: generation in progress}} \\
Ministral-3-8B & Mistral (Mistral AI) & Kaggle GPU & \multicolumn{8}{l}{\pending{extension: generation in progress}} \\
\midrule Total & & & 2{,}368 & 11{,}656 & & & & & & \\
\bottomrule
\end{tabular}

\end{table*}

\begin{table}[t]
\centering\caption{Scope of this analysis and planned extensions.}\label{tab:status}
\scriptsize% Generated by report/build.py -- do not edit.
\begin{tabular}{@{}p{0.31\columnwidth}p{0.33\columnwidth}p{0.24\columnwidth}@{}}
\toprule
Item & Analysed here & Next \\
\midrule
Gemini 3.5 Flash-Lite, generation & 8 regions analysed (1{,}259 recipes) & subset set by the free-tier daily request cap \\
Qwen3.5-4B, scale contrast (RQ4) & planned extension & extension: generation in progress \\
Ministral-3-8B, third open-weight family & planned extension & extension: generation in progress \\
Best-of-$N$ steering & GPT-6-Luna, Gemini 3.5 Flash-Lite and Qwen3.5-9B & extension: Llama-3.1-8B \\
Critique-and-revise & GPT-6-Luna (2 rounds) and Qwen3.5-9B (1 round) & extension: other models \\
Ingredient aliasing & automatic CulinaryDB protocol (Appendix~\ref{app:alias}) & manual curation of rare terms planned \\
Distribution of D1, D4, D5 & packaged and tested locally & public release planned \\
\bottomrule
\end{tabular}

\end{table}

% =====================================================================================================
\section{Methodology}
\subsection{Food pairing and null models}
For a recipe $R$ with $n\geq 2$ ingredients and flavor-compound sets $C_i$~\cite{ahn2011,jain2015},
\begin{equation}
\Ns(R)=\frac{2}{n(n-1)}\sum_{i<j\in R}|C_i\cap C_j| .
\end{equation}
A cuisine's $\bar N_s$ is the mean over its recipes, $\dNs=\bar N_s^{\mathrm{cuisine}}-\bar N_s^{\mathrm{rand}}$, $d=\dNs/\sigma_{\mathrm{rand}}$ and
\begin{equation}
Z=\sqrt{N_{\mathrm{rand}}}\,\frac{\bar N_s^{\mathrm{cuisine}}-\bar N_s^{\mathrm{rand}}}{\sigma_{\mathrm{rand}}}=\sqrt{N_{\mathrm{rand}}}\;d ,
\end{equation}
with $N_{\mathrm{rand}}=\NNRand$ and $\sigma_{\mathrm{rand}}$ the standard deviation of recipe $\Ns$ in the null cuisine~\cite{jain2015,singh2018}. Because $Z$ grows with $\sqrt{N_{\mathrm{rand}}}$, not with the cuisine's size, we always report $d$ and $\dNs$ with it. The four nulls keep the cuisine's own ingredient pool and recipe-size distribution and draw ingredients uniformly (\emph{random}), in proportion to their frequency (\emph{frequency}), within each recipe's category composition (\emph{category}), or both~\cite{singh2018}. For LLM, copy-mutate and genuine-subsample cuisines the null is estimated from \NNullCandidates{} recipes and $Z$ keeps $N_{\mathrm{rand}}=\NNRand$, so it stays on the published scale. Choices the paper leaves open were settled by a calibration grid against the published values: ingredients without molecules stay in recipes (zero-sharing pairs), additives are kept, and Cayenne uses Ahn et al.'s profile.

\subsection{Drivers, fingerprints and composition}
The contribution $\chi\%_i$ is the percentage change in the cuisine's mean $\Ns$ when ingredient $i$ is removed from its recipes, oriented by the sign of the cuisine's own $Z$, among the \NChiTopK{} most frequent non-additive ingredients (the removal variant that best reproduces the published contributors~\cite{singh2018}); LLM and real drivers are compared by the Jaccard overlap of their top five and the Spearman correlation of $\chi\%$. The category fingerprint re-draws only the ingredients of one category in every recipe and records $\Delta N_s^{\mathrm{cat}}$ and $Z^{\mathrm{cat}}$ for each category~\cite{jain2015}; we compare fingerprints by Spearman $\rho$ over shared categories. Prevalence is $P_{ic}=n_{ic}/N_c$ and relative prevalence $p_{ic}=P_{ic}-\mathrm{mean}_{c'\neq c}P_{ic'}$~\cite{jain2018}; the signature set $\Sigma_c$ holds the top-\NSignatureK{} real ingredients by $p_{ic}$, and the stereotype amplification index is $\mathrm{SAI}_c=\sum_{i\in\Sigma_c}P^{\mathrm{LLM}}_{ic}/\sum_{i\in\Sigma_c}P^{\mathrm{real}}_{ic}$ (1 = as in real recipes). We also report the Jensen--Shannon divergence (JSD, base 2) of normalized prevalence and of category composition, the frequency-rank MAE~\cite{tuwani2019}, the diversity ratio (unique ingredients over those of the real cuisine rarefied to the same $N$) and novelty: each recipe's maximum Jaccard similarity to any CulinaryDB recipe, a near-copy being $\geq\NNearCopy$ (genuine recipes are scored leave-one-out).

\subsection{Co-occurrence networks}
For each cuisine we rank ingredient pairs by normalized PMI~\cite{bouma2009}, $\mathrm{NPMI}_{ij}=\log\frac{p_{ij}}{p_ip_j}\big/(-\log p_{ij})$, among pairs found in at least $\max(3, 2\%)$ of its recipes, and compare the top 100 pairs (primary; top 50 as a sensitivity check) with the full real cuisine's by Jaccard similarity. We also extract the disparity-filter backbone of the count-weighted co-occurrence network~\cite{serrano2009} and compare it by Jaccard, by overlap $|A\cap B|/|A|$ (the share of a cuisine's backbone edges that are also real backbone edges) and by size. The reference is \NNetDraws{} genuine held-out samples per region of the LLM cuisine's own size, each compared with the remaining recipes so that no recipe is on both sides; LLM values get a one-sided empirical $p$ (less real-like than genuine), with BH across regions within each model.

\subsection{Aliasing}
We re-implemented the CulinaryDB protocol~\cite{singh2018}: lower-casing, removal of punctuation and numbers, singularization, English and culinary stop-words, longest-first $n$-gram matching against entity names and synonyms (including CulinaryDB's synonym notation), partial matches flagged, and tie-breaks learned from expert lines. On the held-out \NAliasTestShare\% of CulinaryDB's raw phrases (\NAliasTestLines{} unique lines) it returns the expert's entity for \NAliasAccLinesPct{} of unique lines and \NAliasAccRowsPct{} of occurrences, with \NAliasCoveragePct{} coverage (Appendix~\ref{app:alias}). LLM lines are aliased the same way, so real and LLM recipes are mapped identically; unmatched lines are dropped, as in CulinaryDB.

\subsection{Baselines and references}
\emph{Genuine equal-$N$ reference.} \NRefSubsamples{} random \NLLMN-recipe subsamples of every real cuisine (fingerprints on \NRefFpSubsamples) are scored by the same function as the LLM cuisines. For food pairing we drew \NEqReps{} subsamples per region at each $N$ from \NEqGridMin{} to \NEqGridMax{} in steps of \NEqGridStep, and re-ran Fig.~4 with every cuisine cut to \NMatchedN{} pairing recipes (\NMatchedReps{} subsamples each).
\emph{Copy-mutate.} Tuwani et al.'s Algorithm~1~\cite{tuwani2019} with Jain and Bagler's ranked fitness (empirical frequency)~\cite{jain2018}: a seed pool of \NCMSeedPool{} ingredients and $m/\phi$ recipes, \NCMMutations{} fitness-biased mutations per copy, the region's real ingredients as the nature basket and its size distribution for the seeds; \NCMInstances{} instances per region, scored on random \NLLMN-recipe subsamples. Random fitness is the secondary variant.

\subsection{Statistics and pre-registration}\label{sec:stats}
Each model $\times$ region cuisine pools both conditions ($N\leq\NRecipesPerRegionModel$); per-condition cuisines ($N\leq\NRecipesPerCondition$) serve RQ4. $\dNs$ has a percentile bootstrap CI over \NBoot{} resamples of recipes. The pre-registered \emph{sign error} is a CI that excludes 0 with the sign opposite to the real cuisine's $Z$. Region is the unit of every test: Fisher's exact test (H1), Wilcoxon signed-rank tests (H2b, H2c, H4, copy-mutate), and Benjamini--Hochberg FDR at $q=\NQ$~\cite{benjamini1995} within each test across models (and across regions for region-level calls). Hypotheses and metrics were fixed in the implementation plan before full-scale LLM data were analysed. One amendment was made after calibrating on genuine samples, when only a pipeline-test summary of one model had been printed (Section~\ref{sec:sizedep}, Appendix~\ref{app:prereg}): each LLM cuisine's own-pool $d$ is placed in the distribution of $d$ over genuine subsamples of the same region and size, with an empirical two-sided $p$ and BH across regions; it is \emph{too uniform} above and \emph{too contrasting} below that distribution.

% =====================================================================================================
\section{Food-Pairing Analysis of the Real Cuisines}
\subsection{Replication gate}
Our pipeline reproduced the published food-pairing $Z$ (Table~\ref{tab:repl}, Fig.~\ref{fig:F2}): \NReplRandomSigns/\NReplRandomN{} signs against the random null, with Spearman $\rho$ of \NReplRhoRange{} across the four nulls, meeting every pre-registered criterion. As in the published Fig.~4, $Z$ against the frequency-preserving nulls is far smaller than against the random null (Fig.~\ref{fig:F2b}, Appendix~\ref{app:extra}), so popularity explains most of the pattern. The published top contributor is in our top three for \NContribTopOne{} regions (top-3 overlap \NContribTopThree); Korea's top contributor is \NKoreaTopOne{} ($\chi\%=\NKoreaTopOneChi$; published \NKoreaTopOnePub). Contrasting cuisines come out more contrasting than published (Korea $Z=\NZKoreaOurs$ vs $\NZKoreaPub$), most likely because today's FlavorDB differs from the 2018 version.

\begin{table}[t]
\centering\caption{Replication of Singh and Bagler's Fig.~4 (configuration \texttt{v1\_keep\_empty}).}\label{tab:repl}
\footnotesize% Generated by report/build.py -- do not edit.
\begin{tabular}{@{}lccl@{}}
\toprule
Null model & Signs & Spearman $\rho$ & Pre-registered criterion \\
\midrule
Random & 22/22 & 0.980 & all signs, $\rho\geq0.90$ \\
Ingredient frequency & 21/22 & 0.980 & $\rho\geq0.70$ \\
Ingredient category & 22/22 & 0.942 & $\rho\geq0.90$ \\
Frequency + category & 21/22 & 0.979 & $\rho\geq0.70$ \\
\midrule \multicolumn{4}{@{}p{0.97\columnwidth}@{}}{Contributions ($\chi\%$): published top-1 in our top-3 for 22/22 regions; top-3 overlap 63/66} \\
\bottomrule
\end{tabular}

\end{table}

\begin{figure*}[t]
\centering\figfile{F2_replication_v1_keep_empty.pdf}{0.92\textwidth}
\figcode{2}\caption{Does our pipeline reproduce the published food-pairing $Z$? Yes: for each null model our $Z$ ($N_{\mathrm{rand}}=\NNRand$) tracks the value read from Singh and Bagler's Fig.~4, with \NReplRandomSigns/\NReplRandomN{} random-null signs and $\rho=\NReplRandomRho$ (blue: published uniform, vermillion: published contrasting cuisines).}\label{fig:F2}
\end{figure*}

\subsection{Food pairing depends on sample size}\label{sec:sizedep}
Before scoring any LLM cuisine we checked the pre-registered sign-error rule on genuine recipes at the LLMs' size. It fired on \NFPRatePct{} of \NRefSamplesTotal{} genuine \NLLMN-recipe subsamples, and on \NFPCanadaPct{} and \NFPChinaPct{} of the Canadian and Chinese ones (Table~\ref{tab:regions}). The cause is the own-pool null (Fig.~\ref{fig:F2d}): a small sample loses the cuisine's long tail of rare ingredients, which share few compounds with the rest, so the null's expected $\Ns$ rises while the sample's mean $\Ns$ does not. For USA, genuine \NLLMN-recipe samples have a median own-pool $\dNs$ of \NUsaOwnPool{} (pool of \NUsaPool{} ingredients) but \NUsaFullPool{} against the full cuisine's pool, the same as the full cuisine's \NUsaFullDeltaNs.

Because \NContrastingAmongSmallest{} of the \NContrastingCount{} contrasting cuisines are among the \NSmallestK{} smallest regions, we re-ran Fig.~4 at matched size (Fig.~\ref{fig:F2c}). The published contrasting cuisines are robust: \NContrastingRobustText{} were contrasting in every \NMatchedN-recipe subsample, and \NUniformRobust{} of the \NUniformCount{} uniform cuisines stayed positive in at least 95\% of subsamples. \NUniformFlipped{} were positive in only \NUniformFlippedShares{} of subsamples, and \NUniformMiddle{} in \NUniformMiddleRange. We read this as a refinement of how $Z$ should be compared across samples, not as a challenge to the published findings: since $Z=\sqrt{N_{\mathrm{rand}}}\,d$ with a size-dependent $d$, comparisons are fair at equal $N$ or against a fixed pool. Our primary RQ1 inference therefore uses genuine samples of the same size; the pre-registered rule is still reported with its false-positive rate.

\begin{figure*}[t]
\centering\figfile{F2d_pool_effect.pdf}{0.92\textwidth}
\figcode{2d}\caption{Why does food pairing depend on the number of recipes? Small samples lose rare, low-sharing ingredients, so the own-pool null's mean rises and $\dNs$ falls (vermillion); against the full cuisine's pool (blue) $\dNs$ stays at the full cuisine's value (dashed). Median over \NEqReps{} genuine subsamples per size.}\label{fig:F2d}
\end{figure*}

\begin{figure}[t]
\centering\figfile{F2c_matched_size_signs.pdf}{\columnwidth}
\figcode{2c}\caption{Does each cuisine's food-pairing sign survive at equal size? For the contrasting cuisines, yes (\NContrastingRobustText{} negative in every \NMatchedN-recipe subsample); \NUniformFlipped{} lose their positive sign. Dot: median $d$ of \NMatchedReps{} subsamples; line: 95\% range; diamond: full cuisine.}\label{fig:F2c}
\end{figure}

\begin{table}[t]
\centering\caption{Per-region robustness of the published sign: full-cuisine $Z$ (ours and published), and how often genuine samples match the sign, trigger the sign-error rule ($N=\NLLMN$) and stay positive ($N=\NMatchedN$). $^\dagger$\,contrasting.}\label{tab:regions}
\scriptsize\setlength{\tabcolsep}{4pt}% Generated by report/build.py -- do not edit.
\begin{tabular}{@{}lrrrrrr@{}}
\toprule
& & \multicolumn{2}{c}{$Z$ vs random null} & \multicolumn{2}{c}{160 recipes} & 290 \\
\cmidrule(lr){3-4}\cmidrule(lr){5-6}\cmidrule(l){7-7}
Region & $N$ & ours & publ. & \makecell{sign\\match} & \makecell{rule\\fires} & \makecell{share\\$d>0$} \\
\midrule
GRC & 931 & 287.9 & 241 & 100\% & 0.0\% & 100\% \\
AFR & 649 & 273.3 & 262 & 100\% & 0.0\% & 100\% \\
ITA & 7{,}466 & 249.3 & 263 & 100\% & 0.0\% & 100\% \\
CBN & 1{,}091 & 229.2 & 250 & 100\% & 0.0\% & 100\% \\
ESP & 813 & 220.6 & 229 & 100\% & 0.0\% & 100\% \\
ME & 993 & 206.1 & 160 & 97\% & 0.1\% & 100\% \\
USA & 16{,}030 & 200.4 & 198 & 52\% & 6.0\% & 86\% \\
INSC & 3{,}979 & 164.0 & 172 & 64\% & 6.0\% & 94\% \\
ANZ & 494 & 117.4 & 110 & 100\% & 0.0\% & 100\% \\
MEX & 3{,}129 & 114.6 & 140 & 72\% & 2.8\% & 92\% \\
THA & 655 & 103.1 & 74 & 87\% & 1.1\% & 98\% \\
FRA & 2{,}684 & 94.4 & 92 & 14\% & 40.8\% & 28\% \\
SAM & 308 & 75.5 & 96 & 94\% & 0.0\% & 100\% \\
SEA & 610 & 43.7 & 48 & 31\% & 17.3\% & 72\% \\
CAN & 1{,}104 & 33.3 & 45 & 3\% & 54.7\% & 16\% \\
CHN & 940 & 29.5 & 59 & 5\% & 61.3\% & 22\% \\
EE$^\dagger$ & 565 & \ensuremath{-}51.0 & \ensuremath{-}31 & 99\% & 0.0\% & 0\% \\
KOR$^\dagger$ & 296 & \ensuremath{-}56.7 & \ensuremath{-}32 & 100\% & 0.0\% & 0\% \\
BRI$^\dagger$ & 1{,}064 & \ensuremath{-}67.6 & \ensuremath{-}48 & 100\% & 0.0\% & 0\% \\
DACH$^\dagger$ & 485 & \ensuremath{-}87.0 & \ensuremath{-}55 & 100\% & 0.0\% & 0\% \\
JPN$^\dagger$ & 572 & \ensuremath{-}123.9 & \ensuremath{-}111 & 100\% & 0.0\% & 0\% \\
SCND$^\dagger$ & 404 & \ensuremath{-}155.0 & \ensuremath{-}148 & 100\% & 0.0\% & 0\% \\
\bottomrule
\end{tabular}

\end{table}

% =====================================================================================================
\section{Fingerprint-Guided Generation}\label{sec:steer}
\subsection{Fidelity scorer}
The scorer adds three components, each standardized against genuine recipes of the cuisine, with \NWeighting{} weights (chosen on validation data by the mean AUC over the three negative types below): (a) the temperature-scaled~\cite{guo2017} log-probability of the cuisine from a \NClfName{} on ingredient indicators and category counts, chosen among linear models (validation macro-F1 \NSvmValFOne{} vs \NLogregValFOneRange{} for logistic regression) because its per-ingredient contributions drive the feedback; (b) the mean shrinkage-smoothed PMI of the recipe's ingredient pairs within the cuisine ($\kappa=\NKappa$); (c) $\Ns$ consistency, $1-|2F_c(\Ns(R))-1|$ with $F_c$ the cuisine's empirical CDF. Splits are \NSplitTrain/\NSplitVal/\NSplitTest{} per region.

The frozen \NClfName{} reached test macro-F1 \NClfFOne, top-1 \NClfTopOne{} and top-3 accuracy \NClfTopThree{} over \NRegions{} cuisines (gradient boosting: validation macro-F1 \NGbdtValFOne), from \NBestRegionF{} down to \NWorstRegionF; temperature scaling cut the expected calibration error from \NEceRaw{} to \NEceCal{} (Table~\ref{tab:scorer}). On the test split the composite separates genuine recipes from other cuisines' (AUC \NAucOtherComposite), copy-mutate (\NAucCMComposite) and random-null recipes (\NAucNullComposite). The components are complementary: the classifier separates cuisines (\NAucOtherClf) but hardly copy-mutate recipes (\NAucCMClf), co-occurrence PMI separates copy-mutate and null recipes (\NAucCMCooc, \NAucNullCooc), and recipe-level $\Ns$ consistency is near chance against every negative type (\NAucNsRange): flavor sharing is a cuisine-level, not a recipe-level, signal.

\begin{table}[t]
\centering\caption{Fidelity scorer on the CulinaryDB test split.}\label{tab:scorer}
\footnotesize\setlength{\tabcolsep}{3pt}% Generated by report/build.py -- do not edit.
\begin{tabular}{@{}lcccc@{}}
\toprule
Macro-F1 / top-1 / top-3 & \multicolumn{4}{l}{0.392 / 0.507 / 0.741} \\
ECE raw $\rightarrow$ scaled & \multicolumn{4}{l}{0.343 $\rightarrow$ 0.060} \\
GBDT comparator, val.\ F1 & \multicolumn{4}{l}{0.422} \\
Independent classifier, F1 & \multicolumn{4}{l}{0.402} \\
\midrule
AUC: genuine vs & Comp. & Classif. & PMI & $N_s$ \\
\midrule
other cuisines & 0.869 & 0.898 & 0.780 & 0.559 \\
copy-mutate & 0.891 & 0.560 & 0.953 & 0.499 \\
random null & 0.943 & 0.626 & 0.989 & 0.537 \\
\bottomrule
\end{tabular}

\end{table}

\subsection{Steering methods}
\emph{Best-of-$N$.} The pre-registered design drew a fresh pool of \NSteerPrePool{} plain-prompt candidates per region for \NSteerPreModels{} on \NSteerPreRegions{} regions, with $N\in\{\NSteerPreN\}$ and \NSteerPreSelected{} selections. For compute reasons we ran a documented deviation: reranking the \emph{already generated} recipes, a pool of up to \NBonPoolMax{} per model $\times$ region (both conditions), every region a model covers, $N\in\{\NBonNValues\}$, up to \NBonSelectedMax{} recipes selected per region and $N$; candidates that are near-copies of a CulinaryDB recipe or duplicate a selected ingredient set are skipped, and ablations drop one component at $N=\NBonNMax$.
\emph{Critique-and-revise} (Self-Refine~\cite{madaan2023}) revises an unsteered random selection of recipes per region (drawn like best-of-$N$ at $N=1$): the feedback names up to \NFeedbackItems{} atypical ingredients (classifier contribution below the cuisine's 5th percentile), up to \NFeedbackItems{} missing signature ingredients that pair well (positive PMI) and whether $\Ns$ lies outside the cuisine's interquartile range, and asks for a revision that keeps the dish (at most \NRoundsMax{} rounds).
\emph{Evaluation} never uses the scorer: an independent gradient-boosting classifier trained on the scorer's held-out split (macro-F1 \NEvalClfFOne) gives the share of recipes assigned to the claimed cuisine; we add prevalence JSD, SAI, diversity, novelty, cuisine-level sign and fingerprint, and, for revisions, the Jaccard similarity of each revised recipe's ingredients to its original. Each cuisine is scored with its recipes in a canonical order, so the seeded nulls give identical values for identical recipe sets: round 0 (the recipes that were revised, read from each revision's metadata) equals best-of-$N$ at $N=1$, apart from the few selected recipes with fewer than two mapped ingredients, which were not sent for revision.

% =====================================================================================================
\section{Results}
\subsection{Food pairing (RQ1)}
Against the full real cuisines (Figs.~\ref{fig:F3},~\ref{fig:F4}), LLM cuisines matched the published sign in \NSignAgreeRange{} of regions (genuine \NLLMN-recipe samples: \NRefSignMatchPct{} on average), and the pre-registered rule flagged \NSignErrText{} sign errors, against \NSignErrExpectedText{} expected from sampling noise alone (Fig.~\ref{fig:F5}). Sign errors did not cluster in the contrasting regions: H1 was supported for \NHOneFisherSupport{} (Fisher, BH-adjusted $p\geq\NHOneFisherPadjMin$). In an exploratory comparison, \NHOneExpDirection{} (median shift in $d$ of \NHOneExpShiftRange); for small samples the pool effect predicts a shift toward negative $d$.

At equal $N$ (Table~\ref{tab:verdicts}), \NEqTypicalTotal{} of \NEqCuisinesTotal{} model--region cuisines were typical of genuine samples, \NEqUniformTotal{} too uniform and \NEqContrastingTotal{} too contrasting; per model, \NEqPerModelText. The test of a shift toward uniform pairing was significant for \NHOneEqWilcoxSupport{} (\NHOneEqWilcoxPmin), so H1 is not supported at equal $N$ either. The direction did depend on region type: LLM $d$ lay above the genuine median in \NDirText{} contrasting vs uniform regions, i.e.\ \NDirConsistent{} (Fisher, raw $p$ \NHOneEqFisherPRange, BH-adjusted \NHOneEqFisherPadjRange). Moreover, \NUnanimousText. \NDirInterpretation

\begin{figure}[t]
\centering\figfile{F3_pairing_dotplot.pdf}{\columnwidth}
\figcode{3}\caption{Do LLM cuisines reproduce each region's food-pairing sign? Only partly: sign agreement with the full real cuisine was \NSignAgreeRange{} of regions, but genuine \NLLMN-recipe samples also miss the sign in some uniform regions (Section~\ref{sec:sizedep}). Grey bands: contrasting regions; symlog axis.}\label{fig:F3}
\end{figure}

\begin{figure*}[t]
\centering\figfile{F4_effect_sizes.pdf}{0.8\textwidth}
\figcode{4}\caption{How large is each model's food-pairing effect? Effect size $d$ with 95\% bootstrap CI per region; the black tick is the full real cuisine, which is not an equal-size reference (Table~\ref{tab:verdicts} gives the equal-$N$ comparison).}\label{fig:F4}
\end{figure*}

\begin{figure*}[t]
\centering\figfile{F5_sign_matrix.pdf}{0.85\textwidth}
\figcode{5}\caption{Where do models get the food-pairing sign wrong under the pre-registered rule? Orange: CI-significant sign error; green: significant and correct; grey: not significant; white dot: region not covered by the model. Many uniform-region errors also occur in genuine samples of the same size (Table~\ref{tab:regions}).}\label{fig:F5}
\end{figure*}

\begin{table}[t]
\centering\caption{Equal-$N$ food pairing: genuine own-pool $d$ at $N=\NLLMN$ and each model's $d$; shaded with an arrow if outside the genuine distribution (BH, $q=\NQ$): $\uparrow$ too uniform, $\downarrow$ too contrasting. $^\dagger$\,contrasting; ``--'' region not covered.}\label{tab:verdicts}
\scriptsize\setlength{\tabcolsep}{3pt}% Generated by report/build.py -- do not edit.
\begin{tabular}{@{}llrrrr@{}}
\toprule
Region & \makecell{Genuine $d$: median\\{[}95\% range{]}} & \makecell{GPT-6-\\Luna} & \makecell{Gemini 3.5\\Flash-Lite} & \makecell{Qwen3.5-\\9B} & \makecell{Llama-3.1-\\8B} \\
\midrule
GRC & 0.64 [0.38, 0.92] & \cellcolor{TooContrasting!25}0.29\,$\downarrow$ & \textcolor{gray}{--} & \cellcolor{TooContrasting!25}0.04\,$\downarrow$ & \cellcolor{TooContrasting!25}0.11\,$\downarrow$ \\
AFR & 0.41 [0.16, 0.65] & 0.48 & 0.40 & 0.50 & \cellcolor{TooContrasting!25}\ensuremath{-}0.31\,$\downarrow$ \\
ITA & 0.37 [0.14, 0.67] & \cellcolor{TooContrasting!25}\ensuremath{-}0.12\,$\downarrow$ & \textcolor{gray}{--} & 0.32 & 0.32 \\
CBN & 0.36 [0.15, 0.61] & \cellcolor{TooContrasting!25}\ensuremath{-}0.01\,$\downarrow$ & \cellcolor{TooContrasting!25}0.03\,$\downarrow$ & \cellcolor{TooContrasting!25}0.01\,$\downarrow$ & 0.10 \\
ESP & 0.51 [0.27, 0.70] & \cellcolor{TooContrasting!25}0.16\,$\downarrow$ & \textcolor{gray}{--} & 0.68 & 0.73 \\
ME & 0.24 [\ensuremath{-}0.02, 0.48] & 0.19 & \textcolor{gray}{--} & 0.37 & 0.00 \\
USA & 0.00 [\ensuremath{-}0.20, 0.19] & \cellcolor{TooContrasting!25}\ensuremath{-}0.48\,$\downarrow$ & \textcolor{gray}{--} & \ensuremath{-}0.02 & \cellcolor{TooUniform!30}0.30\,$\uparrow$ \\
INSC & 0.04 [\ensuremath{-}0.19, 0.29] & \cellcolor{TooContrasting!25}\ensuremath{-}0.48\,$\downarrow$ & \textcolor{gray}{--} & \cellcolor{TooContrasting!25}\ensuremath{-}0.45\,$\downarrow$ & \ensuremath{-}0.14 \\
ANZ & 0.21 [0.00, 0.37] & \cellcolor{TooContrasting!25}\ensuremath{-}0.45\,$\downarrow$ & \cellcolor{TooContrasting!25}\ensuremath{-}0.27\,$\downarrow$ & 0.07 & 0.25 \\
MEX & 0.07 [\ensuremath{-}0.15, 0.30] & 0.05 & \textcolor{gray}{--} & \ensuremath{-}0.13 & \cellcolor{TooUniform!30}0.84\,$\uparrow$ \\
THA & 0.15 [\ensuremath{-}0.07, 0.39] & \cellcolor{TooContrasting!25}\ensuremath{-}0.22\,$\downarrow$ & \textcolor{gray}{--} & \ensuremath{-}0.13 & 0.11 \\
FRA & \ensuremath{-}0.13 [\ensuremath{-}0.34, 0.11] & \cellcolor{TooUniform!30}0.34\,$\uparrow$ & \textcolor{gray}{--} & \cellcolor{TooUniform!30}0.39\,$\uparrow$ & \cellcolor{TooUniform!30}0.45\,$\uparrow$ \\
SAM & 0.15 [\ensuremath{-}0.05, 0.34] & \ensuremath{-}0.08 & \textcolor{gray}{--} & \cellcolor{TooUniform!30}0.63\,$\uparrow$ & \cellcolor{TooUniform!30}1.08\,$\uparrow$ \\
SEA & \ensuremath{-}0.05 [\ensuremath{-}0.25, 0.16] & \cellcolor{TooContrasting!25}\ensuremath{-}0.49\,$\downarrow$ & \textcolor{gray}{--} & \ensuremath{-}0.07 & \cellcolor{TooUniform!30}0.28\,$\uparrow$ \\
CAN & \ensuremath{-}0.18 [\ensuremath{-}0.37, 0.01] & \ensuremath{-}0.32 & \ensuremath{-}0.28 & \ensuremath{-}0.14 & \ensuremath{-}0.22 \\
CHN & \ensuremath{-}0.16 [\ensuremath{-}0.35, 0.04] & \cellcolor{TooContrasting!25}\ensuremath{-}0.73\,$\downarrow$ & \cellcolor{TooContrasting!25}\ensuremath{-}0.63\,$\downarrow$ & \ensuremath{-}0.29 & \cellcolor{TooUniform!30}0.24\,$\uparrow$ \\
EE$^\dagger$ & \ensuremath{-}0.26 [\ensuremath{-}0.46, \ensuremath{-}0.06] & \ensuremath{-}0.35 & \cellcolor{TooUniform!30}0.03\,$\uparrow$ & \ensuremath{-}0.12 & \ensuremath{-}0.36 \\
KOR$^\dagger$ & \ensuremath{-}0.22 [\ensuremath{-}0.38, \ensuremath{-}0.07] & \ensuremath{-}0.15 & \textcolor{gray}{--} & \ensuremath{-}0.07 & \cellcolor{TooUniform!30}0.11\,$\uparrow$ \\
BRI$^\dagger$ & \ensuremath{-}0.42 [\ensuremath{-}0.62, \ensuremath{-}0.23] & \cellcolor{TooUniform!30}0.25\,$\uparrow$ & \cellcolor{TooUniform!30}0.09\,$\uparrow$ & \cellcolor{TooUniform!30}0.00\,$\uparrow$ & \cellcolor{TooUniform!30}0.00\,$\uparrow$ \\
DACH$^\dagger$ & \ensuremath{-}0.43 [\ensuremath{-}0.64, \ensuremath{-}0.26] & \ensuremath{-}0.33 & \cellcolor{TooUniform!30}\ensuremath{-}0.21\,$\uparrow$ & \ensuremath{-}0.47 & \ensuremath{-}0.26 \\
JPN$^\dagger$ & \ensuremath{-}0.51 [\ensuremath{-}0.69, \ensuremath{-}0.34] & \cellcolor{TooContrasting!25}\ensuremath{-}0.92\,$\downarrow$ & \textcolor{gray}{--} & \ensuremath{-}0.52 & \ensuremath{-}0.47 \\
SCND$^\dagger$ & \ensuremath{-}0.56 [\ensuremath{-}0.74, \ensuremath{-}0.40] & \ensuremath{-}0.63 & \textcolor{gray}{--} & \ensuremath{-}0.42 & \cellcolor{TooUniform!30}\ensuremath{-}0.24\,$\uparrow$ \\
\bottomrule
\end{tabular}

\end{table}

\subsection{Flavor drivers (RQ2a)}
The top-5 $\chi\%$ contributors of LLM cuisines overlapped those of the real cuisines with median Jaccard \NDrvJaccardRange{}, and the rank correlation of $\chi\%$ had median \NDrvRhoRange{}, against \NDrvRefJaccard{} and \NDrvRefRho{} for genuine \NLLMN-recipe samples (\NDrvRefReps{} per region); per model, \NDrvPerModelText. \NDrvNegText{} \NDrvKoreaText{} (real top five: \NKoreaRealTopFive)\iffigure{F7_drivers.pdf}{; Fig.~\ref{fig:F7}}.

\IfFileExists{../figures/F7_drivers.pdf}{%
\begin{figure}[t]
\centering\figfile{F7_drivers.pdf}{\columnwidth}
\figcode{7}\caption{Do LLM recipes rely on the same flavor drivers? Korea's real top contributor, \NKoreaTopOne, keeps its role for some models but not all; across regions the overlap of top-5 drivers (median \NDrvJaccardRange) stays below that of genuine samples (\NDrvRefJaccard).}\label{fig:F7}
\end{figure}}{}

\subsection{Composition, stereotypes and diversity (RQ2b)}
LLM cuisines used ingredients in different proportions than real ones: median prevalence JSD \NPrevJSDRange{} vs \NRefPrevJSD{} for genuine samples, and frequency-rank MAE \NFreqMAERange{} vs \NRefFreqMAE{} (Table~\ref{tab:scorecard}, Fig.~\ref{fig:F12}). Contrary to H2b they did not over-use signature ingredients on average: median SAI was \NSAIRange{} (genuine 95\% band \NRefSAIBandLo--\NRefSAIBandHi), and $\log\mathrm{SAI}>0$ was supported for \NHTwoBSupport{} (BH-adjusted $p\geq\NHTwoBPadjMin$), although single regions deviate in both directions (Fig.~\ref{fig:F8}). What the models lose is breadth: the diversity ratio was \NDivRange{}, a reduction supported for \NHTwoBDivSupport{} (BH-adjusted $p\leq\NHTwoBDivPadjMax$)\iffigure{F9_diversity.pdf}{; Fig.~\ref{fig:F9}}. \NNearCopyRangePct{} of LLM recipes were near-copies of a CulinaryDB recipe, similar to the \NRefNearCopyPct{} of genuine recipes that have a near-duplicate elsewhere in CulinaryDB\iffigure{F10_sizes.pdf}{; recipe sizes are compared in Fig.~\ref{fig:F10}}.

\begin{figure}[t]
\centering\figfile{F8_sai.pdf}{\columnwidth}
\figcode{8}\caption{Do LLMs over-use each cuisine's signature ingredients? Not on average: median SAI is \NSAIRange, close to the genuine median of \NRefSAI, but single regions fall outside the range of genuine \NLLMN-recipe samples (grey) in both directions.}\label{fig:F8}
\end{figure}

\IfFileExists{../figures/F9_diversity.pdf}{%
\begin{figure*}[t]
\centering\figfile{F9_diversity.pdf}{0.8\textwidth}
\figcode{9}\caption{Do LLM cuisines use as many different ingredients as real ones? No. (a) For Italy, the LLM rank-frequency curves follow the real cuisine's common ingredients but fall below it in the tail. (b) Across regions, the diversity ratio (unique ingredients relative to genuine samples of the same size, which have ratio 1) is \NDivRange{}, below genuine samples and below copy-mutate.}\label{fig:F9}
\end{figure*}}{}

\IfFileExists{../figures/F10_sizes.pdf}{%
\begin{figure}[t]
\centering\figfile{F10_sizes.pdf}{\columnwidth}
\figcode{10}\caption{Are LLM recipes the same size as real ones? Recipe-size distributions (ingredients per recipe) of the real cuisines and each model.}\label{fig:F10}
\end{figure}}{}

\subsection{Category fingerprints and co-occurrence networks (RQ2c)}
Fingerprints agreed with the real cuisine's at median $\rho$ \NFpRhoRange, against \NRefFpRho{} for genuine samples\iffigure{F6_fingerprints.pdf}{ (Fig.~\ref{fig:F6})}. H2c was supported for \NHTwoCSupport{} (\NHTwoCSummary), and the number of regions below the genuine resampling distribution after BH was \NHTwoCBelowText.

The co-occurrence networks tell the same story at the level of ingredient pairs (Table~\ref{tab:networks}\iffigure{F11_networks.pdf}{, Fig.~\ref{fig:F11}}). The top-100 NPMI pairs of LLM cuisines overlapped the real cuisine's with median Jaccard \NNetHundredRange, against \NNetHundredRef{} for genuine held-out samples of the same size; the overlap was below the genuine range in \NNetHundredBelowText{} regions (top-50 sensitivity: \NNetFiftyRange{} vs \NNetFiftyRef). LLM backbones were \NNetSparser{} genuine ones (median \NNetBbEdgesRange{} vs \NNetBbEdgesRef{} edges), yet \NNetBbOverlapRange{} of their edges were real backbone edges (genuine: \NNetBbOverlapRef). LLMs thus reproduce a generic skeleton of strong pairings but miss each cuisine's distinctive pairings\iffigure{F11b_showcase_pairs.pdf}{ (Fig.~\ref{fig:F11b} shows Korea and Italy)}.

\IfFileExists{../figures/F6_fingerprints.pdf}{%
\begin{figure*}[t]
\centering\figfile{F6_fingerprints.pdf}{0.85\textwidth}
\figcode{6}\caption{Do LLM cuisines keep each cuisine's category fingerprint? Only partly. (a, b) $\Delta N_s^{\mathrm{cat}}$ per category for Korea and Italy (real: full cuisine; dash: category in fewer than 10 of that source's recipes). (c) Across regions, the category contributions $Z^{\mathrm{cat}}$ of LLM cuisines correlate with the real ones at median Spearman $\rho$ \NFpRhoRange, well below genuine samples of the same size (\NRefFpRho).}\label{fig:F6}
\end{figure*}}{}

\begin{table}[t]
\centering\caption{Co-occurrence networks vs the full real cuisine (medians over regions; in parentheses, regions below the genuine range after BH).}\label{tab:networks}
\scriptsize\setlength{\tabcolsep}{2.5pt}% Generated by report/build.py -- do not edit.
\begin{tabular}{@{}lccccc@{}}
\toprule
Metric & \makecell{Genuine\\@160} & \makecell{GPT-6-\\Luna} & \makecell{Gemini 3.5\\Flash-Lite} & \makecell{Qwen3.5-\\9B} & \makecell{Llama-3.1-\\8B} \\
\midrule
\makecell[l]{NPMI top-100,\\Jaccard (primary)} & 0.30 & \makecell{0.11\\(20/22)} & \makecell{0.11\\(8/8)} & \makecell{0.13\\(19/22)} & \makecell{0.13\\(19/22)} \\
\makecell[l]{NPMI top-50,\\Jaccard} & 0.22 & \makecell{0.07\\(20/22)} & \makecell{0.07\\(7/8)} & \makecell{0.07\\(19/22)} & \makecell{0.09\\(16/22)} \\
\makecell[l]{Backbone,\\Jaccard} & 0.29 & \makecell{0.09\\(22/22)} & \makecell{0.09\\(8/8)} & \makecell{0.12\\(20/22)} & \makecell{0.08\\(22/22)} \\
\makecell[l]{Backbone overlap\\$|A\cap B|/|A|$} & 0.98 & \makecell{0.92\\(0/22)} & \makecell{0.85\\(3/8)} & \makecell{0.89\\(14/22)} & \makecell{0.94\\(6/22)} \\
\makecell[l]{Backbone\\edges} & 160 & \makecell{69\\(21/22)} & \makecell{73\\(8/8)} & \makecell{104\\(18/22)} & \makecell{67\\(22/22)} \\
\bottomrule
\end{tabular}

\end{table}

\IfFileExists{../figures/F11_networks.pdf}{%
\begin{figure*}[t]
\centering\figfile{F11_networks.pdf}{0.85\textwidth}
\figcode{11}\caption{Do LLM cuisines combine the same distinctive ingredient pairs as the real cuisines? Mostly not: per region, the overlap of each model's top-100 NPMI pairs with the real cuisine's (median \NNetHundredRange) sits below the range of genuine held-out samples of the same size (median \NNetHundredRef).}\label{fig:F11}
\end{figure*}}{}

\IfFileExists{../figures/F11b_showcase_pairs.pdf}{%
\begin{figure*}[t]
\centering\figfile{F11b_showcase_pairs.pdf}{0.85\textwidth}
\figcode{11b}\caption{Which of a cuisine's distinctive pairings do the models reproduce? The real cuisine's twelve most distinctive (highest-NPMI) pairs for Korea and Italy, the showcase regions fixed in advance; each cell gives the pair's NPMI in that source, and ``absent'' marks pairs found in fewer than $\max(3, 2\%)$ of that model's recipes.}\label{fig:F11b}
\end{figure*}}{}

\subsection{Scorecard and copy-mutate baseline}
At equal $N$, LLM cuisines were closer to the real cuisines than copy-mutate in prevalence JSD for \NCMJSDSupport, in fingerprint $\rho$ for \NCMFpSupport{} and in $|\log\mathrm{SAI}|$ for \NCMSaiSupport{} (Table~\ref{tab:rq}; Fig.~\ref{fig:F13} in Appendix~\ref{app:extra}). Copy-mutate reproduces the rank-frequency curve (MAE \NCMFreqMAE) but not which ingredients a cuisine uses (prevalence JSD \NCMPrevJSD, SAI \NCMSAI, fingerprint $\rho$ \NCMFpRho). On pairing, $|d-d_{\mathrm{full}}|$ was smaller than copy-mutate's for \NCMDgapSupport. Table~\ref{tab:scorecard} and Fig.~\ref{fig:F12} summarize every metric.

\begin{table*}[t]
\centering\caption{Flavor-fidelity scorecard: medians over regions for each model; genuine \NLLMN-recipe samples and copy-mutate (ranked fitness, $N=\NLLMN$) as references (genuine and copy-mutate: medians over regions and replicates; sign rates and near-copy rate of models: means).}\label{tab:scorecard}
\footnotesize% Generated by report/build.py -- do not edit.
\begin{tabular}{@{}lrrrrrr@{}}
\toprule
Metric & \makecell{Genuine\\@160} & \makecell{Copy-mutate\\@160} & \makecell{GPT-6-\\Luna} & \makecell{Gemini 3.5\\Flash-Lite} & \makecell{Qwen3.5-\\9B} & \makecell{Llama-3.1-\\8B} \\
\midrule
Regions (recipes) & -- & -- & 22 (3{,}517) & 8 (1{,}259) & 22 (3{,}477) & 22 (3{,}403) \\
\midrule \multicolumn{7}{@{}l}{\emph{Food pairing (RQ1)}} \\
Sign agrees with real cuisine & 78\% & 46\% & 50\% & 38\% & 68\% & 82\% \\
Pre-registered sign-error rate & 8.6\% & 40\% & 36\% (8/22) & 38\% (3/8) & 14\% (3/22) & 9\% (2/22) \\
Equal-$N$: too uniform / typical / too contrasting & -- & -- & 2 / 9 / 11 & 3 / 2 / 3 & 3 / 16 / 3 & 9 / 11 / 2 \\
Median shift $d-\tilde d_{\mathrm{genuine}}$ & -- & -- & \ensuremath{-}0.29 & \ensuremath{-}0.05 & \ensuremath{-}0.01 & 0.11 \\
\midrule \multicolumn{7}{@{}l}{\emph{Composition, stereotypes and fingerprints (RQ2)}} \\
Prevalence JSD & 0.031 & 0.624 & 0.195 & 0.184 & 0.161 & 0.203 \\
Category-composition JSD & 0.0024 & -- & 0.048 & 0.034 & 0.037 & 0.041 \\
Frequency-rank MAE & 0.0075 & 0.046 & 0.051 & 0.039 & 0.031 & 0.040 \\
SAI & 1.00 & 0.17 & 1.04 & 0.98 & 0.97 & 1.01 \\
Fingerprint $\rho$ & 0.75 & 0.02 & 0.33 & 0.37 & 0.30 & 0.28 \\
Diversity ratio & -- & -- & 0.46 & 0.53 & 0.60 & 0.48 \\
\midrule \multicolumn{7}{@{}l}{\emph{Novelty}} \\
Near-copy rate (max.\ Jaccard $\geq$0.8) & 12.5\% & 0.0\% & 17.6\% & 16.7\% & 9.7\% & 13.4\% \\
Mean max.\ Jaccard to CulinaryDB & -- & -- & 0.68 & 0.65 & 0.61 & 0.64 \\
\bottomrule
\end{tabular}

\end{table*}

\begin{figure*}[t]
\centering\figfile{F12_scorecard.pdf}{0.72\textwidth}
\figcode{12}\caption{How faithful is each source to the real cuisines? Genuine samples set the ceiling on every metric; the LLMs sit between them and copy-mutate on composition and fingerprints. Medians over the regions each source covers.}\label{fig:F12}
\end{figure*}

\subsection{Prompt, vendor and scale (RQ4)}
The persona lowered prevalence JSD for \NPersonaJSDSupport{} (median change \NPersonaJSDRange), raised fingerprint $\rho$ for \NPersonaFpSupport, and changed $|d-d_{\mathrm{full}}|$ for \NPersonaDgapSupport{} and $|\log\mathrm{SAI}|$ for \NPersonaSaiSupport{} (Fig.~\ref{fig:F14}). Averaged per region, the proprietary models (\NPropModels) differed from the open ones (\NOpenModels) by \NPropDgap{} in $|d-d_{\mathrm{full}}|$ ($p_{\mathrm{BH}}=\NPropDgapP$), \NPropJSD{} in prevalence JSD ($p_{\mathrm{BH}}=\NPropJSDP$), \NPropFp{} in fingerprint $\rho$ ($p_{\mathrm{BH}}=\NPropFpP$) and \NPropSai{} in $|\log\mathrm{SAI}|$ ($p_{\mathrm{BH}}=\NPropSaiP$): \NPropVerdict. \NPropCaveat{} Scale: \NScaleText.

\begin{figure*}[t]
\centering\figfile{F14_persona_vs_plain.pdf}{0.8\textwidth}
\figcode{14}\caption{Does the home-cook persona change fidelity? It lowers prevalence JSD for \NPersonaJSDSupport; effects on fingerprints and SAI are mixed. One point per region (persona minus plain); bar: median.}\label{fig:F14}
\end{figure*}

\begin{table*}[t]
\centering\caption{Pre-registered tests (region as unit). Cells: effect (BH-adjusted $p$), bold if $p_{\mathrm{BH}}<\NQ$. Effects are medians over regions unless a count is shown; Appendix~\ref{app:tests} lists every statistic.}\label{tab:rq}
\footnotesize% Generated by report/build.py -- do not edit.
\begin{tabular}{@{}lcccc@{}}
\toprule
Test (effect) & GPT-6-Luna & Gemini 3.5 Flash-Lite & Qwen3.5-9B & Llama-3.1-8B \\
\midrule
H1: sign errors, contrasting vs uniform (Fisher) & 1/6 vs 7/16 (1.00) & 0/3 vs 3/5 (1.00) & 0/6 vs 3/16 (1.00) & 0/6 vs 2/16 (1.00) \\
H1: opposite-sign regions, size-aware $p$ (BH) & 8/22 & 3/8 & 2/22 & 2/22 \\
H1 (expl.): $d-d_{\mathrm{full}}>0$ & \ensuremath{-}0.493 (1.00) & \ensuremath{-}0.427 (1.00) & \ensuremath{-}0.264 (1.00) & \ensuremath{-}0.101 (1.00) \\
H1 equal-$N$: $d-\tilde d_{\mathrm{gen}}>0$ & \ensuremath{-}0.287 (0.99) & \ensuremath{-}0.054 (0.91) & \ensuremath{-}0.012 (0.91) & 0.114 (0.26) \\
H1 equal-$N$: shift up, contrasting vs uniform & 3/6 vs 2/16 (0.20) & 3/3 vs 0/5 (0.071) & 4/6 vs 6/16 (0.23) & 5/6 vs 8/16 (0.23) \\
\midrule
H2b: $\log\mathrm{SAI}>0$ (median SAI) & 1.042 (0.64) & 0.984 (0.88) & 0.975 (0.88) & 1.009 (0.88) \\
H2b: $\log$ diversity ratio $<0$ (median) & \textbf{0.459 (\ensuremath{<}0.001)} & \textbf{0.532 (0.004)} & \textbf{0.597 (\ensuremath{<}0.001)} & \textbf{0.477 (\ensuremath{<}0.001)} \\
\midrule
H2c: $\rho-\tilde\rho_{\mathrm{gen}}<0$ & \textbf{\ensuremath{-}0.416 (\ensuremath{<}0.001)} & \textbf{\ensuremath{-}0.328 (0.004)} & \textbf{\ensuremath{-}0.384 (\ensuremath{<}0.001)} & \textbf{\ensuremath{-}0.480 (\ensuremath{<}0.001)} \\
H2c: regions below genuine range (BH) & 14/22 & 4/8 & 20/22 & 19/22 \\
\midrule
H4 persona$-$plain: $|d-d_{\mathrm{full}}|$ & \ensuremath{-}0.151 (0.18) & \ensuremath{-}0.065 (0.38) & \ensuremath{-}0.162 (0.18) & \ensuremath{-}0.082 (0.28) \\
H4 persona$-$plain: prevalence JSD & \textbf{\ensuremath{-}0.076 (\ensuremath{<}0.001)} & \textbf{\ensuremath{-}0.092 (0.016)} & \textbf{\ensuremath{-}0.041 (\ensuremath{<}0.001)} & \textbf{\ensuremath{-}0.034 (\ensuremath{<}0.001)} \\
H4 persona$-$plain: fingerprint $\rho$ & 0.104 (0.33) & 0.504 (0.33) & \textbf{0.114 (0.005)} & 0.006 (0.71) \\
H4 persona$-$plain: $|\log\mathrm{SAI}|$ & \ensuremath{-}0.081 (0.26) & 0.002 (0.95) & \ensuremath{-}0.064 (0.15) & \ensuremath{-}0.048 (0.15) \\
\midrule
LLM$-$copy-mutate: $|d-d_{\mathrm{full}}|$ & \ensuremath{-}0.082 (0.46) & \ensuremath{-}0.021 (0.46) & \textbf{\ensuremath{-}0.148 (0.005)} & \textbf{\ensuremath{-}0.165 (0.005)} \\
LLM$-$copy-mutate: prevalence JSD & \textbf{\ensuremath{-}0.441 (\ensuremath{<}0.001)} & \textbf{\ensuremath{-}0.444 (0.008)} & \textbf{\ensuremath{-}0.452 (\ensuremath{<}0.001)} & \textbf{\ensuremath{-}0.432 (\ensuremath{<}0.001)} \\
LLM$-$copy-mutate: fingerprint $\rho$ & \textbf{0.262 (\ensuremath{<}0.001)} & \textbf{0.423 (0.039)} & \textbf{0.300 (0.001)} & \textbf{0.229 (0.001)} \\
LLM$-$copy-mutate: $|\log\mathrm{SAI}|$ & \textbf{\ensuremath{-}1.840 (\ensuremath{<}0.001)} & \textbf{\ensuremath{-}1.922 (0.008)} & \textbf{\ensuremath{-}1.860 (\ensuremath{<}0.001)} & \textbf{\ensuremath{-}1.778 (\ensuremath{<}0.001)} \\
\midrule
H4 proprietary$-$open: $|d-d_{\mathrm{full}}|$ & \multicolumn{4}{c}{0.095 (0.074)} \\
H4 proprietary$-$open: prevalence JSD & \multicolumn{4}{c}{\textbf{0.016 (0.007)}} \\
H4 proprietary$-$open: fingerprint $\rho$ & \multicolumn{4}{c}{0.036 (0.50)} \\
H4 proprietary$-$open: $|\log\mathrm{SAI}|$ & \multicolumn{4}{c}{0.020 (0.82)} \\
\bottomrule
\end{tabular}

\end{table*}

\subsection{Steering (RQ3)}
\textbf{Best-of-$N$.} The independent classifier's share of recipes assigned to the claimed cuisine \NBonClfMonotone{} (Fig.~\ref{fig:F15}; Table~\ref{tab:bon} in Appendix~\ref{app:extra}): \NBonClfText. Diversity \NBonDivMonotone{} (\NBonDivText), and SAI rose from \NBonSaiText, so reranking makes recipes more typical partly by leaning on signature ingredients; novelty fell (mean maximum Jaccard \NBonNoveltyText). Cuisine-level metrics did not improve: sign agreement went from \NBonSignText, fingerprint $\rho$ from \NBonFpText{} and prevalence JSD from \NBonJSDText. With at most \NBonSelectedMax{} recipes per region these cuisine-level statistics are too noisy to show an effect. The near-copy guard held (rate \NBonNearCopy{} at every $N$). In the ablation at $N=\NBonNMax$ (Fig.~\ref{fig:F15b}, Appendix~\ref{app:extra}), averaged over models, \NBonAblationText.

\textbf{Critique-and-revise.} We ran \NReviseRoundsText{} (Table~\ref{tab:revise}, Fig.~\ref{fig:F15c}). The independent classifier's share went from \NReviseClfText; sign agreement from \NReviseSignText; fingerprint $\rho$ from \NReviseFpText; prevalence JSD from \NReviseJSDText. Revised recipes kept \NReviseIdentityPctRange{} of their original ingredients (Jaccard). With \NBonSelectedMax{} recipes per region, the cuisine-level changes are suggestive only.

\begin{table}[t]
\centering\caption{Critique-and-revise (round 0: the unsteered recipes that were revised; held-out metrics; means over regions).}\label{tab:revise}
\scriptsize\setlength{\tabcolsep}{2.5pt}% Generated by report/build.py -- do not edit.
\begin{tabular}{@{}lcccccc@{}}
\toprule
Model & Round & \makecell{Indep.\\classif.} & \makecell{Sign\\agree.} & \makecell{Finger-\\print $\rho$} & \makecell{Prev.\\JSD} & \makecell{Identity\\(Jaccard)} \\
\midrule
GPT-6-Luna & 0 & 0.45 & 0.41 & 0.25 & 0.264 & -- \\
 & 1 & 0.53 & 0.59 & 0.31 & 0.247 & 0.81 \\
 & 2 & 0.58 & 0.59 & 0.35 & 0.241 & 0.77 \\
Qwen3.5-9B & 0 & 0.47 & 0.36 & 0.38 & 0.259 & -- \\
 & 1 & 0.55 & 0.68 & 0.41 & 0.245 & 0.80 \\
\bottomrule
\end{tabular}

\end{table}

\begin{figure*}[t]
\centering\figfile{F15_best_of_n.pdf}{0.92\textwidth}
\figcode{15}\caption{Does best-of-$N$ reranking with the fidelity scorer make recipes more faithful? It raises the independent classifier's agreement with the claimed cuisine, but also SAI, while diversity and novelty fall and prevalence JSD does not improve. Means over regions with 95\% bootstrap CIs.}\label{fig:F15}
\end{figure*}

\begin{figure}[t]
\centering\figfile{F15c_critique_revise.pdf}{\columnwidth}
\figcode{15c}\caption{Does fingerprint feedback improve recipes when the model revises them? The independent agreement and prevalence JSD improve with each round while revised recipes keep most of their original ingredients; cuisine-level changes rest on few recipes per region.}\label{fig:F15c}
\end{figure}

% =====================================================================================================
\section{Discussion}
\textbf{What LLMs get right.} Judged against genuine samples of the same size, most LLM cuisines have a typical food-pairing effect, and none of the models shifts cuisines toward uniform pairing as H1 predicted. They use each cuisine's signature ingredients at about the real rate, and they are far closer to real cuisines than a copy-mutate model that knows each cuisine's ingredient popularity. Most of the strong pairings they produce are pairings the real cuisine also relies on: their backbone edges are mostly real backbone edges.

\textbf{What they miss.} The losses are in breadth and specificity. LLM cuisines draw on about half the ingredient repertoire of a real cuisine of the same size, they share few of the distinctive pairings that NPMI picks out, and their category fingerprints agree with the real ones far less than resamples do. The picture is a generic skeleton of popular, compatible ingredients: plausible recipes that recognisably belong to a cuisine but flatten its long tail. This is the chemistry-level counterpart of the textual finding that LLMs do not ground cultural adaptation in salient ingredients~\cite{carichon2026}. Signature ingredients are therefore not where the stereotyping shows; the narrowing of everything around them is.

\textbf{Models and prompts.} The proprietary models were not more faithful than the open-weight ones, and the models differ in direction as much as in degree (equal-$N$ verdicts, Table~\ref{tab:verdicts}). A home-cook persona, which varies the place in every call, brings ingredient prevalence closer to the real cuisine. That is a cheap, prompt-level improvement.

\textbf{Steering.} Reranking with a scorer that includes a cuisine classifier makes recipes more recognizable to an independent classifier, but it does so by selecting typical recipes, which raises signature-ingredient use and lowers diversity and novelty, a typicality trade-off the held-out metrics expose. Critique-and-revise, which names specific atypical and missing ingredients, improved the independent agreement and prevalence JSD while keeping most of each dish. Cuisine-level pairing and fingerprints need many more steered recipes per region before they can show an effect.

\textbf{Implications for the framework.} The size dependence of the own-pool null matters beyond LLMs: any comparison of $Z$ between cuisines or samples of different sizes mixes the pairing effect with the size of the ingredient pool. Comparing at equal $N$ or against a fixed pool keeps the published contrasting cuisines intact and gives a fair yardstick for generated cuisines.

% =====================================================================================================
\section{Limitations}
(i) \emph{One run per model.} Each model $\times$ region cuisine comes from one set of \NReplicates{} calls per condition; LLM outputs are stochastic and hosted models change, so we log model ids, versions and seeds. (ii) \emph{Small cuisines.} \NRecipesPerRegionModel{} recipes per region (\NRecipesPerCondition{} per condition) make pairing and fingerprints noisy; equal-$N$ references account for sample size, but not for the LLMs' narrower ingredient pools, and power is limited. Steering used at most \NBonSelectedMax{} recipes per region. (iii) \emph{Steering design.} Best-of-$N$ deviated from the pre-registered design. The independent classifier is a different model class trained on different recipes, but it learns from the same kind of data as the scorer, so part of its agreement may reflect shared inductive bias. (iv) \emph{Aliasing.} \NAliasMatchedRange{} of LLM lines were aliased and unmatched lines are dropped; curation of the remaining terms awaits review. (v) \emph{Flavor data.} Today's FlavorDB is not the 2018 version and \NMissingProfiles{} profiles are reconstructed; rankings of contributors replicate, magnitudes only approximately. (vi) \emph{Coverage.} Gemini 3.5 Flash-Lite covers \NGeminiRegions{} of \NRegions{} regions and the scale contrast is missing (Table~\ref{tab:status}). The proprietary models use provider defaults while open models use $T=\NTemperature$. (vii) \emph{Amendment.} The equal-$N$ comparison was added after pre-registration (before the LLM analysis); we report the pre-registered rule as well. (viii) Recipes were neither cooked nor tasted; flavor sharing is a proxy for flavor.

% =====================================================================================================
\section{Conclusion}
After reproducing Singh and Bagler's food-pairing results, we found that their own-pool null is size-dependent; comparing at equal size keeps \NContrastingRobustText{} contrasting cuisines robust and gives a fair yardstick for LLM cuisines. Against that yardstick, the cuisines of \NModelsDone{} LLMs \NFindingPairing{} and \NFindingSAI, but \NFindingFingerprint, they drew on only \NDivRangePct{} of the real cuisines' ingredient diversity, and they reproduced a generic skeleton of strong pairings rather than each cuisine's distinctive ones. \NFindingCM{} Fingerprint-guided steering made recipes more recognizable to an independent classifier; critique-and-revise did so while keeping most of each dish, whereas best-of-$N$ traded diversity and novelty for typicality. LLMCuisine-22, FlavorAudit and FlavorPrint make the audit repeatable for any recipe collection.

\section*{Acknowledgments}
We thank Prof.\ Ganesh Bagler and the course TA. CulinaryDB and FlavorDB are resources of the Complex Systems Lab, IIIT-Delhi.

\bibliographystyle{IEEEtran}
\bibliography{refs}

% =====================================================================================================
\appendices

\section{Deliverables and traceability}\label{app:trace}
\textbf{D1 LLMCuisine-22} (\texttt{generation/}, \texttt{data/release/LLMCuisine-22/}): recipes in CulinaryDB's file format plus model, condition and call id, the full generation log, a datasheet and a coverage report. \textbf{D2 Audit} (\texttt{analysis/}, \texttt{tables/}, \texttt{figures/}): replication, equal-$N$ reference, copy-mutate baseline, scorecard, networks and pre-registered tests, all regenerated by scripts. \textbf{D3 Steering} (\texttt{steering/}, \texttt{flavoraudit.scoring}): scorer, best-of-$N$ and critique-and-revise. \textbf{D4 FlavorAudit} (\texttt{src/flavoraudit/}): a pip-installable MIT library with \NLibModules{} modules, \NLibTests{} test functions, a command-line interface and a documentation site of \NDocsPages{} pages; data are downloaded on first use. \textbf{D5 FlavorPrint} (\texttt{app/}): a Streamlit app with \NAppTabs{} tabs (audit a pasted recipe against a claimed cuisine; results dashboard; methods and licences), packaged for Hugging Face Spaces. Table~\ref{tab:trace} maps the proposal to what was delivered.
\begin{table*}[t]
\centering\caption{Traceability from the submitted proposal to the delivered work.}\label{tab:trace}
\footnotesize% Generated by report/build.py -- do not edit.
\begin{tabular}{@{}p{0.11\textwidth}p{0.32\textwidth}p{0.2\textwidth}p{0.3\textwidth}@{}}
\toprule
Deliverable & Proposed (Oct.\ 2026) & Where & Status at this build \\
\midrule
D1 LLMCuisine-22 & $\geq$10{,}000 recipes, 22 regions, $\geq$3 LLMs, plain + persona prompts, CulinaryDB format, datasheet & \texttt{generation/}, \texttt{data/release/}, datasheet & 11{,}656 recipes from four LLMs so far; targets met \\
D2 Audit & $N_s$, $Z$ vs four nulls with CIs, contributions, fingerprints, SAI, diversity, novelty; copy-mutate baseline; scorecard & \texttt{analysis/}, \texttt{tables/}, \texttt{figures/} & replication, equal-$N$ reference, scorecard, pre-registered tests, $\chi\%$ drivers, co-occurrence networks done \\
D3 Steering & Fidelity scorer (classifier + co-occurrence + $N_s$) steering best-of-$N$ and critique-and-revise & \texttt{steering/}, \texttt{flavoraudit.scoring} & scorer trained and validated; best-of-$N$ done (deviated design); critique-and-revise done \\
D4 FlavorAudit & pip-installable, tested, documented library (proposed as ``CuisineKit'') & \texttt{src/flavoraudit/}, \texttt{tests/}, MkDocs site & 14 modules, 99 tests, documentation built; PyPI release awaits the team's approval \\
D5 FlavorPrint & Streamlit app on Hugging Face Spaces: audit a recipe; dashboard & \texttt{app/} & Audit, Dashboard and About tabs built and tested locally; deployment awaits the team's approval \\
\bottomrule
\end{tabular}

\end{table*}

\section{Aliasing evaluation}\label{app:alias}
Unique lower-cased raw phrases were split by a stable hash into training (fit and validation) and test (\NAliasTestShare\%). Variants were compared on validation lines with what was learned on the fit lines; the chosen variant (tie-break \NAliasChosen) was refit on all training lines and evaluated once on the test lines (Table~\ref{tab:alias}). Accuracy by source ranged \NAliasSourceRange; the weakest source was \NAliasWorstSource.
\begin{table}[t]
\centering\caption{Aliaser variants (validation) and the held-out test result. $^\ast$\,chosen.}\label{tab:alias}
\scriptsize\setlength{\tabcolsep}{3pt}% Generated by report/build.py -- do not edit.
\begin{tabular}{@{}lp{0.44\columnwidth}rrr@{}}
\toprule
& Variant & \makecell{Unique\\lines} & \makecell{Occur-\\rences} & Coverage \\
\midrule
A & Protocol, all CulinaryDB entities, no partial matches & 89.1\% & 92.2\% & 97.7\% \\
B & + vocabulary of entities used by the curators & 89.2\% & 92.3\% & 97.6\% \\
C & + partial matches & 91.4\% & 93.4\% & 100.0\% \\
D$^\ast$ & + tie-break preferences learned from expert lines & 97.8\% & 98.8\% & 100.0\% \\
D & + extended culinary stop-words & 97.4\% & 98.6\% & 99.9\% \\
E & + mined synonyms (support $\geq$5, precision $\geq$0.9) & 97.8\% & 98.8\% & 100.0\% \\
F & + mined synonyms (support $\geq$10, precision $\geq$0.95) & 97.8\% & 98.8\% & 100.0\% \\
\midrule
Test & chosen variant, refit on all training lines & 97.9\% & 99.0\% & 100.0\% \\
\midrule
\multicolumn{5}{@{}p{0.97\columnwidth}}{Test accuracy by source: Allrecipes 98.9\%, Epicurious 97.6\%, Food Network 97.5\%, Tarla Dalal 90.0\%} \\
\bottomrule
\end{tabular}

\end{table}

\section{Pre-registration and amendments}\label{app:prereg}
The research questions, hypotheses, metrics and tests (Section~I) were written into the implementation plan before any full-scale LLM output was analysed; later changes are dated in its decision log. Those affecting the analysis: (1) decoding was set to a uniform temperature \NTemperature{} and top-$p$ \NTopP{} for open-weight models and provider defaults for APIs, before any data existed; (2) the prompts were frozen after the pilot (prompt set v1); (3) the RQ4 scale pair became Qwen3.5-9B vs Qwen3.5-4B when the larger OpenAI model was dropped from the roster to save cost, before any full run; (4) the equal-$N$ food-pairing comparison was added to RQ1 after the genuine-sample calibration of Section~\ref{sec:sizedep}, when only a pipeline-test summary of one model had been printed; (5) the scorer's weights were selected on validation data against all three negative types; (6) best-of-$N$ was run on the already generated recipes instead of a fresh pool (Section~\ref{sec:steer}).

\section{All test statistics}\label{app:tests}
\begin{table*}[t]
\centering\caption{Every pre-registered and exploratory test (\texttt{tables/rq\_tests\_full.csv}); $p_{\mathrm{BH}}$ within each test across models.}\label{tab:rqfull}
\tiny\setlength{\tabcolsep}{3pt}% Generated by report/build.py -- do not edit.
\begin{tabular}{@{}llp{0.34\textwidth}rrrrr@{}}
\toprule
Hyp. & Model & Test & $n$ & Statistic & Effect & $p$ & $p_{\mathrm{BH}}$ \\
\midrule
H1 & real@160 & sign-error rule false-positive rate & 22{,}000 & -- & 0.086 & -- & -- \\
H1 & Gemini 3.5 Flash-Lite & Fisher exact, errors contrasting > uniform & 8 & 0.00 & 0/3 vs 3/5 & 1.00 & 1.00 \\
H1 & Gemini 3.5 Flash-Lite & sign errors, size-aware p with BH across regions & 8 & -- & 3.000 & -- & -- \\
H1-exploratory & Gemini 3.5 Flash-Lite & Wilcoxon d\_LLM - d\_real > 0 (push toward positive pairing) & 8 & 6.00 & \ensuremath{-}0.427 & 0.96 & 1.00 \\
H1 & Llama-3.1-8B & Fisher exact, errors contrasting > uniform & 22 & 0.00 & 0/6 vs 2/16 & 1.00 & 1.00 \\
H1 & Llama-3.1-8B & sign errors, size-aware p with BH across regions & 22 & -- & 2.000 & -- & -- \\
H1-exploratory & Llama-3.1-8B & Wilcoxon d\_LLM - d\_real > 0 (push toward positive pairing) & 22 & 90.00 & \ensuremath{-}0.101 & 0.88 & 1.00 \\
H1 & GPT-6-Luna & Fisher exact, errors contrasting > uniform & 22 & 0.26 & 1/6 vs 7/16 & 0.96 & 1.00 \\
H1 & GPT-6-Luna & sign errors, size-aware p with BH across regions & 22 & -- & 8.000 & -- & -- \\
H1-exploratory & GPT-6-Luna & Wilcoxon d\_LLM - d\_real > 0 (push toward positive pairing) & 22 & 14.00 & \ensuremath{-}0.493 & 1.00 & 1.00 \\
H1 & Qwen3.5-9B & Fisher exact, errors contrasting > uniform & 22 & 0.00 & 0/6 vs 3/16 & 1.00 & 1.00 \\
H1 & Qwen3.5-9B & sign errors, size-aware p with BH across regions & 22 & -- & 2.000 & -- & -- \\
H1-exploratory & Qwen3.5-9B & Wilcoxon d\_LLM - d\_real > 0 (push toward positive pairing) & 22 & 37.00 & \ensuremath{-}0.264 & 1.00 & 1.00 \\
H1-equalN & Gemini 3.5 Flash-Lite & Wilcoxon d\_LLM - median d\_genuine(same N) > 0 & 8 & 15.00 & \ensuremath{-}0.054 & 0.68 & 0.91 \\
H1-equalN & Gemini 3.5 Flash-Lite & Fisher: shift toward uniform, contrasting vs uniform regions & 8 & $\infty$ & 3/3 vs 0/5 & 0.018 & 0.071 \\
H1-equalN & Llama-3.1-8B & Wilcoxon d\_LLM - median d\_genuine(same N) > 0 & 22 & 174.00 & 0.114 & 0.064 & 0.26 \\
H1-equalN & Llama-3.1-8B & Fisher: shift toward uniform, contrasting vs uniform regions & 22 & 5.00 & 5/6 vs 8/16 & 0.18 & 0.23 \\
H1-equalN & GPT-6-Luna & Wilcoxon d\_LLM - median d\_genuine(same N) > 0 & 22 & 55.00 & \ensuremath{-}0.287 & 0.99 & 0.99 \\
H1-equalN & GPT-6-Luna & Fisher: shift toward uniform, contrasting vs uniform regions & 22 & 7.00 & 3/6 vs 2/16 & 0.10 & 0.20 \\
H1-equalN & Qwen3.5-9B & Wilcoxon d\_LLM - median d\_genuine(same N) > 0 & 22 & 128.00 & \ensuremath{-}0.012 & 0.49 & 0.91 \\
H1-equalN & Qwen3.5-9B & Fisher: shift toward uniform, contrasting vs uniform regions & 22 & 3.33 & 4/6 vs 6/16 & 0.23 & 0.23 \\
H2b & Gemini 3.5 Flash-Lite & Wilcoxon log SAI > 0 & 8 & 10.00 & 0.984 & 0.88 & 0.88 \\
H2b-diversity & Gemini 3.5 Flash-Lite & Wilcoxon log diversity ratio < 0 & 8 & 0.00 & 0.532 & 0.004 & 0.004 \\
H2b & Llama-3.1-8B & Wilcoxon log SAI > 0 & 22 & 115.00 & 1.009 & 0.65 & 0.88 \\
H2b-diversity & Llama-3.1-8B & Wilcoxon log diversity ratio < 0 & 22 & 0.00 & 0.477 & \ensuremath{<}0.001 & \ensuremath{<}0.001 \\
H2b & GPT-6-Luna & Wilcoxon log SAI > 0 & 22 & 158.00 & 1.042 & 0.16 & 0.64 \\
H2b-diversity & GPT-6-Luna & Wilcoxon log diversity ratio < 0 & 22 & 0.00 & 0.459 & \ensuremath{<}0.001 & \ensuremath{<}0.001 \\
H2b & Qwen3.5-9B & Wilcoxon log SAI > 0 & 22 & 102.00 & 0.975 & 0.79 & 0.88 \\
H2b-diversity & Qwen3.5-9B & Wilcoxon log diversity ratio < 0 & 22 & 0.00 & 0.597 & \ensuremath{<}0.001 & \ensuremath{<}0.001 \\
H2b & real@160 & median SAI of genuine samples (expected 1) & 22{,}000 & -- & 1.000 & -- & -- \\
H2c & Gemini 3.5 Flash-Lite & Wilcoxon rho\_LLM - median rho\_real@160 < 0 & 8 & 0.00 & \ensuremath{-}0.328 & 0.004 & 0.004 \\
H2c & Llama-3.1-8B & Wilcoxon rho\_LLM - median rho\_real@160 < 0 & 22 & 0.00 & \ensuremath{-}0.480 & \ensuremath{<}0.001 & \ensuremath{<}0.001 \\
H2c & GPT-6-Luna & Wilcoxon rho\_LLM - median rho\_real@160 < 0 & 22 & 4.00 & \ensuremath{-}0.416 & \ensuremath{<}0.001 & \ensuremath{<}0.001 \\
H2c & Qwen3.5-9B & Wilcoxon rho\_LLM - median rho\_real@160 < 0 & 22 & 0.00 & \ensuremath{-}0.384 & \ensuremath{<}0.001 & \ensuremath{<}0.001 \\
H4-persona & Gemini 3.5 Flash-Lite & Wilcoxon persona - plain: abs\_d\_gap & 8 & 11.00 & \ensuremath{-}0.065 & 0.38 & 0.38 \\
H4-persona & Gemini 3.5 Flash-Lite & Wilcoxon persona - plain: prevalence\_jsd & 8 & 1.00 & \ensuremath{-}0.092 & 0.016 & 0.016 \\
H4-persona & Gemini 3.5 Flash-Lite & Wilcoxon persona - plain: fingerprint\_spearman & 8 & 9.00 & 0.504 & 0.25 & 0.33 \\
H4-persona & Gemini 3.5 Flash-Lite & Wilcoxon persona - plain: abs\_log\_sai & 8 & 17.00 & 0.002 & 0.95 & 0.95 \\
H4-persona & Llama-3.1-8B & Wilcoxon persona - plain: abs\_d\_gap & 22 & 87.00 & \ensuremath{-}0.082 & 0.21 & 0.28 \\
H4-persona & Llama-3.1-8B & Wilcoxon persona - plain: prevalence\_jsd & 22 & 15.00 & \ensuremath{-}0.034 & \ensuremath{<}0.001 & \ensuremath{<}0.001 \\
H4-persona & Llama-3.1-8B & Wilcoxon persona - plain: fingerprint\_spearman & 22 & 115.00 & 0.006 & 0.71 & 0.71 \\
H4-persona & Llama-3.1-8B & Wilcoxon persona - plain: abs\_log\_sai & 22 & 63.00 & \ensuremath{-}0.048 & 0.039 & 0.15 \\
H4-persona & GPT-6-Luna & Wilcoxon persona - plain: abs\_d\_gap & 22 & 74.00 & \ensuremath{-}0.151 & 0.092 & 0.18 \\
H4-persona & GPT-6-Luna & Wilcoxon persona - plain: prevalence\_jsd & 22 & 0.00 & \ensuremath{-}0.076 & \ensuremath{<}0.001 & \ensuremath{<}0.001 \\
H4-persona & GPT-6-Luna & Wilcoxon persona - plain: fingerprint\_spearman & 22 & 88.00 & 0.104 & 0.22 & 0.33 \\
H4-persona & GPT-6-Luna & Wilcoxon persona - plain: abs\_log\_sai & 22 & 86.00 & \ensuremath{-}0.081 & 0.20 & 0.26 \\
H4-persona & Qwen3.5-9B & Wilcoxon persona - plain: abs\_d\_gap & 22 & 69.00 & \ensuremath{-}0.162 & 0.063 & 0.18 \\
H4-persona & Qwen3.5-9B & Wilcoxon persona - plain: prevalence\_jsd & 22 & 3.00 & \ensuremath{-}0.041 & \ensuremath{<}0.001 & \ensuremath{<}0.001 \\
H4-persona & Qwen3.5-9B & Wilcoxon persona - plain: fingerprint\_spearman & 22 & 32.00 & 0.114 & 0.001 & 0.005 \\
H4-persona & Qwen3.5-9B & Wilcoxon persona - plain: abs\_log\_sai & 22 & 71.00 & \ensuremath{-}0.064 & 0.074 & 0.15 \\
H4-proprietary & Gemini 3.5 Flash-Lite + GPT-6-Luna & Wilcoxon proprietary - open: abs\_d\_gap & 22 & 71.00 & 0.095 & 0.074 & 0.074 \\
H4-proprietary & Gemini 3.5 Flash-Lite + GPT-6-Luna & Wilcoxon proprietary - open: prevalence\_jsd & 22 & 45.00 & 0.016 & 0.007 & 0.007 \\
H4-proprietary & Gemini 3.5 Flash-Lite + GPT-6-Luna & Wilcoxon proprietary - open: fingerprint\_spearman & 22 & 105.00 & 0.036 & 0.50 & 0.50 \\
H4-proprietary & Gemini 3.5 Flash-Lite + GPT-6-Luna & Wilcoxon proprietary - open: abs\_log\_sai & 22 & 119.00 & 0.020 & 0.82 & 0.82 \\
CM & Gemini 3.5 Flash-Lite & Wilcoxon LLM - CM-R/ranked: abs\_d\_gap & 8 & 12.00 & \ensuremath{-}0.021 & 0.46 & 0.46 \\
CM & Llama-3.1-8B & Wilcoxon LLM - CM-R/ranked: abs\_d\_gap & 22 & 37.00 & \ensuremath{-}0.165 & 0.003 & 0.005 \\
CM & GPT-6-Luna & Wilcoxon LLM - CM-R/ranked: abs\_d\_gap & 22 & 103.00 & \ensuremath{-}0.082 & 0.46 & 0.46 \\
CM & Qwen3.5-9B & Wilcoxon LLM - CM-R/ranked: abs\_d\_gap & 22 & 32.00 & \ensuremath{-}0.148 & 0.001 & 0.005 \\
CM & Gemini 3.5 Flash-Lite & Wilcoxon LLM - CM-R/ranked: prevalence\_jsd & 8 & 0.00 & \ensuremath{-}0.444 & 0.008 & 0.008 \\
CM & Llama-3.1-8B & Wilcoxon LLM - CM-R/ranked: prevalence\_jsd & 22 & 0.00 & \ensuremath{-}0.432 & \ensuremath{<}0.001 & \ensuremath{<}0.001 \\
CM & GPT-6-Luna & Wilcoxon LLM - CM-R/ranked: prevalence\_jsd & 22 & 0.00 & \ensuremath{-}0.441 & \ensuremath{<}0.001 & \ensuremath{<}0.001 \\
CM & Qwen3.5-9B & Wilcoxon LLM - CM-R/ranked: prevalence\_jsd & 22 & 0.00 & \ensuremath{-}0.452 & \ensuremath{<}0.001 & \ensuremath{<}0.001 \\
CM & Gemini 3.5 Flash-Lite & Wilcoxon LLM - CM-R/ranked: fingerprint\_spearman & 8 & 3.00 & 0.423 & 0.039 & 0.039 \\
CM & Llama-3.1-8B & Wilcoxon LLM - CM-R/ranked: fingerprint\_spearman & 22 & 26.00 & 0.229 & \ensuremath{<}0.001 & 0.001 \\
CM & GPT-6-Luna & Wilcoxon LLM - CM-R/ranked: fingerprint\_spearman & 22 & 18.00 & 0.262 & \ensuremath{<}0.001 & \ensuremath{<}0.001 \\
CM & Qwen3.5-9B & Wilcoxon LLM - CM-R/ranked: fingerprint\_spearman & 22 & 30.00 & 0.300 & \ensuremath{<}0.001 & 0.001 \\
CM & Gemini 3.5 Flash-Lite & Wilcoxon LLM - CM-R/ranked: abs\_log\_sai & 8 & 0.00 & \ensuremath{-}1.922 & 0.008 & 0.008 \\
CM & Llama-3.1-8B & Wilcoxon LLM - CM-R/ranked: abs\_log\_sai & 22 & 0.00 & \ensuremath{-}1.778 & \ensuremath{<}0.001 & \ensuremath{<}0.001 \\
CM & GPT-6-Luna & Wilcoxon LLM - CM-R/ranked: abs\_log\_sai & 22 & 0.00 & \ensuremath{-}1.840 & \ensuremath{<}0.001 & \ensuremath{<}0.001 \\
CM & Qwen3.5-9B & Wilcoxon LLM - CM-R/ranked: abs\_log\_sai & 22 & 0.00 & \ensuremath{-}1.860 & \ensuremath{<}0.001 & \ensuremath{<}0.001 \\
\bottomrule
\end{tabular}

\end{table*}

\section{Additional figures and tables}\label{app:extra}
\begin{figure}[t]
\centering\figfile{F2b_fig4_style_v1_keep_empty.pdf}{\columnwidth}
\figcode{2b}\caption{Do our results look like Singh and Bagler's Fig.~4? Yes: filled markers (ours) sit next to the open published ones for all four nulls, and the frequency-preserving nulls bring $Z$ close to zero. Grey bands: contrasting cuisines.}\label{fig:F2b}
\end{figure}

\begin{table}[t]
\centering\caption{Best-of-$N$ reranking (deviated design; held-out metrics only; means over regions).}\label{tab:bon}
\scriptsize\setlength{\tabcolsep}{3pt}% Generated by report/build.py -- do not edit.
\begin{tabular}{@{}lccc@{}}
\toprule
Metric (mean over regions) & \makecell{GPT-6-\\Luna} & \makecell{Gemini 3.5\\Flash-Lite} & \makecell{Qwen3.5-\\9B} \\
\midrule
\multicolumn{4}{@{}l}{\emph{Independent classifier: share in claimed cuisine}} \\
\quad $N=1$ & 0.45 & 0.31 & 0.47 \\
\quad $N=2$ & 0.53 & 0.34 & 0.50 \\
\quad $N=4$ & 0.59 & 0.40 & 0.56 \\
\quad $N=8$ & 0.61 & 0.47 & 0.60 \\
\multicolumn{4}{@{}l}{\emph{Diversity ratio}} \\
\quad $N=1$ & 0.69 & 0.76 & 0.76 \\
\quad $N=2$ & 0.67 & 0.69 & 0.71 \\
\quad $N=4$ & 0.65 & 0.66 & 0.69 \\
\quad $N=8$ & 0.60 & 0.66 & 0.63 \\
\multicolumn{4}{@{}l}{\emph{Prevalence JSD}} \\
\quad $N=1$ & 0.264 & 0.275 & 0.259 \\
\quad $N=2$ & 0.258 & 0.280 & 0.255 \\
\quad $N=4$ & 0.261 & 0.276 & 0.259 \\
\quad $N=8$ & 0.275 & 0.269 & 0.270 \\
\multicolumn{4}{@{}l}{\emph{Fingerprint $\rho$}} \\
\quad $N=1$ & 0.25 & 0.34 & 0.38 \\
\quad $N=2$ & 0.18 & 0.32 & 0.15 \\
\quad $N=4$ & 0.17 & 0.23 & 0.18 \\
\quad $N=8$ & 0.18 & 0.03 & 0.13 \\
\multicolumn{4}{@{}l}{\emph{Sign agreement (share of regions)}} \\
\quad $N=1$ & 0.41 & 0.50 & 0.36 \\
\quad $N=2$ & 0.32 & 0.25 & 0.50 \\
\quad $N=4$ & 0.50 & 0.50 & 0.50 \\
\quad $N=8$ & 0.41 & 0.50 & 0.45 \\
\multicolumn{4}{@{}l}{\emph{Near-copy rate}} \\
\quad $N=1$ & 0.00 & 0.00 & 0.00 \\
\quad $N=2$ & 0.00 & 0.00 & 0.00 \\
\quad $N=4$ & 0.00 & 0.00 & 0.00 \\
\quad $N=8$ & 0.00 & 0.00 & 0.00 \\
\bottomrule
\end{tabular}

\end{table}

\begin{figure*}[t]
\centering\figfile{F13_copy_mutate.pdf}{0.86\textwidth}
\figcode{13}\caption{Do LLMs capture cuisines better than a copy-mutate evolution model? On which ingredients a cuisine uses and on its fingerprint, yes; on the rank-frequency curve, no. One point per region.}\label{fig:F13}
\end{figure*}

\begin{figure}[t]
\centering\figfile{F15b_ablation.pdf}{\columnwidth}
\figcode{15b}\caption{Which scorer component drives the best-of-$N$ gain? The classifier: dropping it lowers the independent agreement most, while dropping $\Ns$ consistency raises agreement but worsens prevalence JSD (ablation at $N=\NBonNMax$).}\label{fig:F15b}
\end{figure}

\section{Frozen prompts}\label{app:prompts}
% Generated by report/build.py -- do not edit.
\noindent\textbf{System:} {\ttfamily\footnotesize You are a helpful cooking assistant.}\par\smallskip
\noindent\textbf{Plain:} {\ttfamily\footnotesize Give me 5 different traditional \{phrase\} recipes.\par \{format\}}\par\smallskip
\noindent\textbf{Persona:} {\ttfamily\footnotesize You are a home cook who grew up in \{place\}, and you still cook the food of your family and region every day. Share 5 different traditional dishes that you actually cook at home.\par \{format\}}\par\smallskip
\noindent\textbf{\{format\}:} {\ttfamily\footnotesize Reply with JSON only (no other text and no markdown), in exactly this structure: \{"recipes": [\{"title": "<dish name>", "ingredients": ["<ingredient line>", "..."], "instructions": ["<step>", "..."]\}]\} The "recipes" list must contain exactly 5 recipes. Write every ingredient line the way it would appear in a recipe, with its quantity, for example "2 cups basmati rice, rinsed". Write the instructions as short steps in cooking order. Use English.}\par\smallskip


\section{Use of AI assistance}
The code, analyses and this report were developed with the help of an AI coding assistant (Claude) under the team's direction and review. The assistant wrote and tested code and drafted text; the team set the research questions, approved every design decision recorded in the decision log, and is responsible for the content. All numbers are produced by the project's scripts, not typed by hand.

\end{document}
